\chapter{稀疏条件与鲁棒重建}
\label{ch:12}

\paragraph*{本章阅读锚点}
本章的锚点是输入不足和输入退化。原始 3DGS 依赖足够视角、可靠 SfM 和干净图像；稀疏鲁棒方法分别补位姿、深度、正则、伪视角或不确定性。


3D Gaussian Splatting 虽然在密集视角、高质量输入条件下展现了卓越的新视角合成能力，但其在真实世界场景中的广泛应用仍面临三大核心挑战。首先，\textbf{稀疏视角}输入（3--9 张图像）下，3DGS 极易对训练视角过拟合，生成充满浮游物和几何畸变的退化表示，这在实际数据采集中极为常见。其次，\textbf{图像退化}（包括运动模糊、散焦模糊、低光照、雾、雨、烟等）破坏了多视角一致性假设，导致 SfM 初始化失败和渲染质量急剧下降。第三，\textbf{相机位姿缺失}使得 3DGS 所依赖的 COLMAP 预处理流程无法运行，而这一前提条件在野外数据采集、长视频和手持拍摄中往往无法满足。此外，\textbf{动态干扰物}（行人、车辆等瞬态物体）、\textbf{水下等特殊环境}的散射介质，以及\textbf{缺乏几何先验}的场景结构模糊性，进一步加剧了重建的不适定性。

本章系统梳理围绕上述挑战展开的 268 篇研究论文，将其组织为以下四个方向：稀疏视角新视角合成、位姿估计与几何先验引导、退化处理与鲁棒重建、主动视觉与扩展应用。通过分析各类技术路线的核心思想、关键创新和适用边界，本章旨在为读者提供在非理想条件下实现鲁棒 3DGS 重建的系统性解决方案。

\begin{figure}[htbp]
\centering
\resizebox{\textwidth}{!}{\begin{tikzpicture}[
  route/.style={rectangle, rounded corners=5pt, draw=#1, fill=#1!10,
    text width=3.2cm, minimum height=0.9cm, align=center, font=\small},
  method/.style={rectangle, rounded corners=3pt, draw=#1, fill=#1!6,
    font=\scriptsize, text width=3.0cm, align=center, minimum height=0.6cm},
  arr/.style={-{Stealth[length=4pt]}, techgray!60},
  >=Stealth
]

  \node[route=techgray, text width=4cm, font=\small\bfseries] (root) at (23.3, 5.6) {稀疏条件与鲁棒重建方法分类};
  \node[route=techblue] (b0) at (7.5, 3.1) {稀疏视角NVS};
  \node[method=techblue] (b0_l0) at (1.7, 0.6) {正则化约束\\{\tiny DropGaussian, CoR-GS, FASR}};
  \node[method=techblue] (b0_l1) at (5.6, 0.6) {扩散先验\\{\tiny ReconX, WaveletGaussian, RI3D}};
  \node[method=techblue] (b0_l2) at (9.4, 0.6) {初始化改进\\{\tiny FSGS, COSMOS, HeroGS}};
  \node[route=techgreen] (b1) at (19.7, 3.1) {位姿估计};
  \node[method=techgreen] (b1_l0) at (13.8, 0.6) {联合位姿优化\\{\tiny GloSplat, JOGS, PCR-GS}};
  \node[method=techgreen] (b1_l1) at (17.9, 0.6) {深度先验\\{\tiny UGOT, Pi-GS, DensifyBeforehand}};
  \node[method=techgreen] (b1_l2) at (21.5, 0.6) {法线约束\\{\tiny UNG-GS, DET-GS, DIP-GS}};
  \node[route=techorange] (b2) at (26.9, 3.1) {退化处理};
  \node[method=techorange] (b2_l0) at (25.1, 0.6) {运动模糊\\{\tiny DeblurGS, CRiM-GS, GeMS}};
  \node[method=techorange] (b2_l1) at (28.7, 0.6) {滚动快门\\{\tiny CRiM-GS, CoMoGaussian}};
  \node[method=techorange] (b2_l2) at (32.1, 0.6) {低光照\\{\tiny Gaussian-DK, LL-Gaussian, Luminance-GS}};
  \node[route=techpurple] (b3) at (35.5, 3.1) {主动视觉};
  \node[method=techpurple] (b3_l0) at (35.5, 0.6) {主动重建\\{\tiny ActiveNeRF-GS, NextBestView-GS}};
  \node[method=techpurple] (b3_l1) at (39.5, 0.6) {最佳视角选择\\{\tiny ExploreGS, Uncertainty-GS}};
  \draw[arr] (root) -- (b0);
  \draw[arr] (root) -- (b1);
  \draw[arr] (root) -- (b2);
  \draw[arr] (root) -- (b3);
  \draw[arr] (b0) -- (b0_l0);
  \draw[arr] (b0) -- (b0_l1);
  \draw[arr] (b0) -- (b0_l2);
  \draw[arr] (b1) -- (b1_l0);
  \draw[arr] (b1) -- (b1_l1);
  \draw[arr] (b1) -- (b1_l2);
  \draw[arr] (b2) -- (b2_l0);
  \draw[arr] (b2) -- (b2_l1);
  \draw[arr] (b2) -- (b2_l2);
  \draw[arr] (b3) -- (b3_l0);
  \draw[arr] (b3) -- (b3_l1);
\end{tikzpicture}}
\caption{稀疏条件与鲁棒重建方法分类}
\label{fig:ch12_sparse_taxonomy}
\end{figure}

\begin{figure}[htbp]
\centering
\resizebox{\textwidth}{!}{\begin{tikzpicture}[
  box/.style={rectangle, rounded corners=4pt, draw=#1, fill=#1!10,
    text width=3.2cm, align=center, font=\small, minimum height=1cm},
  sbox/.style={rectangle, rounded corners=3pt, draw=#1!60, fill=#1!5,
    font=\scriptsize, text width=2.4cm, align=center, minimum height=0.5cm},
  arr/.style={-{Stealth[length=6pt]}, thick, techgray},
  darr/.style={-{Stealth[length=6pt]}, thick, techgray, dashed},
  >=Stealth
]

  \node[box=techblue] (l0) at (9.5, 10.6) {输入少(2-3张)};
  \node[box=techblue] (l1) at (7.2, 8.1) {几何歧义};
  \node[box=techblue] (l2) at (5.9, 5.6) {遮挡严重};
  \node[box=techgreen] (r0) at (12.5, 8.1) {输入多(50-100张)};
  \node[box=techgreen] (r1) at (12.6, 5.6) {几何明确};
  \node[box=techgreen] (r2) at (12.8, 3.1) {充分覆盖};
  \node[sbox=techblue] (ls0) at (1.7, 3.1) {正则化约束};
  \node[sbox=techblue] (ls1) at (5.6, 3.1) {扩散先验};
  \node[sbox=techblue] (ls2) at (9.2, 3.1) {深度估计};
  \node[sbox=techgreen] (rs0) at (12.8, 0.6) {标准3DGS};
  \draw[arr] (l0) -- (l1);
  \draw[arr] (l0) -- (r0);
  \draw[arr] (l1) -- (l2);
  \draw[arr] (l2) -- (ls0);
  \draw[arr] (l2) -- (ls1);
  \draw[arr] (l2) -- (ls2);
  \draw[arr] (r0) -- (r1);
  \draw[arr] (r1) -- (r2);
  \draw[arr] (r2) -- (rs0);
\end{tikzpicture}}
\caption{稀疏视角与稠密视角重建方法对比}
\label{fig:ch12_sparse_comparison}
\end{figure}

% ======================================================================
\section{稀疏视角新视角合成}
\label{sec:12_sparse_nvs}

稀疏视角新视角合成是 3DGS 领域中研究最为密集的子方向之一。当输入视角从数十张锐减至 3--9 张时，3DGS 面临的核心问题是严重的过拟合：高斯基元仅覆盖训练视角可见的区域，对未见区域缺乏约束，导致渲染新视角时出现大面积浮游物、半透明伪影和几何扭曲。现有方法可大致分为四大技术路线：（1）基于正则化与 dropout 的约束方法；（2）基于扩散模型生成的伪视角增强方法；（3）基于初始化改进和致密化的基础设施工具；（4）基于深度先验的几何引导方法。本节按前三条路线展开论述，第四条路线在第~\ref{ss:12_4}~节的深度正则化部分讨论。

\subsection{正则化与几何约束}
\label{ss:12_1}

本节介绍通过约束高斯基元优化行为来抑制稀疏视角过拟合的正则化与 dropout 技术。

\subsubsection{正则化与 Dropout 策略}

本路线介绍通过 dropout 和正则化约束高斯基元优化行为以抑制过拟合的方法。

正则化策略的核心思路是通过约束高斯基元的优化行为来抑制过拟合。其中，\textbf{dropout 类方法}（在训练过程中随机屏蔽部分模型组件以防止过拟合的技术）因其简单有效而成为一条重要研究线。

\paragraph{结构化随机丢弃}

先理解一个关键事实：在稀疏视角下，远离相机的高斯基元往往被前方基元遮挡，其透射率 $T_i$ 极低，导致这些基元几乎收不到有效的梯度反馈，属性（尺度、颜色、不透明度等）得不到充分更新。结果就是整个场景的优化信号几乎全部集中在少数训练视角可见的近处区域，这恰恰是过拟合的温床。

DropGaussian（2025.04）~\cite{2504_00773} 的核心思想受启发于神经网络中的 dropout（一种通过随机丢弃网络单元来防止过拟合的正则化技术）：在每次训练迭代中，随机丢弃一部分高斯基元，使其暂时退出渲染过程。直觉上来说，当前方遮挡基元被丢弃后，原本被遮挡的远处基元突然变得"可见"，从而获得更丰富的梯度信号，有机会更新自身属性。这一操作完全不需要引入外部先验模块，仅通过训练阶段的随机扰动即可显著缓解过拟合。

具体实现上，DropGaussian 在每次迭代时以丢弃率 $r$ 随机选择高斯基元予以移除。例如 $r = 0.1$ 表示丢弃 10\% 的基元。为了维持渲染亮度的一致性，被丢弃基元的贡献通过补偿因子分配给保留基元：保留基元的不透明度乘以 $M(i) = 1/(1-r)$，被丢弃基元的 $M(i) = 0$。这样即使在丢弃过程中，每个像素的总颜色贡献量仍然保持不变，渲染结果不会因丢弃操作而变暗或变亮。

更进一步，DropGaussian 观察到一个重要现象：随着训练推进，过拟合倾向会逐步加剧。因此丢弃率不应该是固定值，而应随训练迭代线性增长。其公式为：

\begin{equation}
    r_t = \lambda \cdot \frac{t}{t_{\text{total}}},
\end{equation}

其中 $\lambda \in (0,1)$ 是缩放因子（默认 $\lambda = 0.2$），$t$ 是当前迭代，$t_{\text{total}}$ 是总迭代数。这一渐进策略在训练初期施加轻微正则化以避免破坏几何建立，后期则加大正则化力度以对抗日益严重的过拟合。消融实验表明，固定丢弃率在 $r = 0.1$ 到 $r = 0.3$ 之间性能逐步下降（过于激进的丢弃反而有害），而渐进策略在所有丢弃率设定下均优于固定策略。此外，随机丢弃优于基于梯度或距离的选择性丢弃，因为后者存在反复丢弃关键重建基元的风险。DropGaussian 仅临时停用基元而非永久移除，被停用的基元在后续迭代中可以重新参与优化。

在定量指标上，DropGaussian 在三个主流基准上均展现了竞争力：LLFF 3-view 设定下达到 PSNR 20.76，Mip-NeRF 360 12-view 设定下取得 PSNR 23.92 / SSIM 0.755 / LPIPS 0.242，Blender 8-view 设定下取得 PSNR 25.42，且不引入任何额外计算开销。

\paragraph{DropoutGS: 基元精简式 dropout}

在 DropGaussian 的基础上，一系列后续工作从不同角度改进了 dropout 策略。DropoutGS（2025.04）~\cite{2504_09491} 发现模型中高斯基元数量越少，过拟合反而越轻，因此提出随机 dropout 正则化（RDR，Random Dropout Regularization）结合边缘引导的分裂策略（ESS，Edge-guided Splitting Strategy），在减少基元数量的同时保留高频细节。

\paragraph{PairDropGS: 一致性正则化 dropout}

PairDropGS（2026.05）~\cite{2605_12399} 从一致性正则化角度重新审视 dropout，构造一对被丢弃的高斯子集，约束其低频渲染结构的一致性，并引入渐进式一致性调度策略以增强训练稳定性。与 DropoutGS 的基元精简思路不同，PairDropGS 侧重于利用被丢弃子集之间的结构一致性来约束优化方向，而非直接减少基元数量。

\paragraph{DropAnSH-GS: 锚点式 dropout 与球谐抑制}

DropAnSH-GS（2026.02）~\cite{2602_20933} 指出已有 dropout 方法忽略了邻近补偿效应（被丢弃的高斯可由邻居补偿）和高阶球谐系数对过拟合的贡献，提出锚点式 dropout 策略并同时丢弃高阶球谐系数。该方法在 DropoutGS 和 PairDropGS 的基础上进一步识别出两种被忽视的过拟合通道，从而提供了更全面的 dropout 正则化。

\paragraph{CoR-GS: 共正则化双辐射场}

除 dropout 外，其他正则化策略同样值得关注。CoR-GS（2024.05）~\cite{2405_12110} 提出共正则化视角：同时训练两个独立的高斯辐射场，利用点不一致性（point disagreement，即两个辐射场对同一三维点位置预测的差异）和渲染不一致性（rendering disagreement，即两个辐射场渲染同一视角时像素颜色的差异）无监督地预测重建质量，并据此剪枝不可靠的高斯。该方法的核心发现是两个辐射场的不一致性与重建精度之间存在稳定的负相关关系，因此可以在不访问真值的情况下识别不准确重建。基于此，CoR-GS 通过共剪枝（co-pruning）移除位置不一致的高斯，并通过伪视角共正则化抑制渲染不一致像素的偏差。

\paragraph{SE-GS: 自集成不确定性正则化}

SE-GS（2024.11）~\cite{2411_00144} 引入自集成策略，通过不确定性感知的扰动训练 $\Delta$-模型和 $\Sigma$-模型，以集成方式增强泛化能力。与 CoR-GS 的双辐射场共训练不同，SE-GS 通过对单一模型的扰动来构建集成，计算成本更为轻量。

\paragraph{FASR: 频率自适应锐度正则化}

FASR（2025.11）~\cite{2511_17918} 将稀疏视角 3DGS 建模为一个泛化问题，提出频率自适应锐度正则化（FASR，Frequency-Adaptive Sharpness Regularization），根据图像局部频率调整正则化强度，在抑制浮游物的同时保留高频细节。

\paragraph{D\texorpdfstring{$^2$}{2}GS: 过拟合-欠拟合双模式治理}

D$^2$GS（2025.10）~\cite{2510_08566} 识别出稀疏视角下过拟合和欠拟合并存的两种失败模式，分别通过深度-密度引导的 dropout 和距离感知的保真度增强来应对。具体而言，近相机区域的高斯基元密度过高导致过拟合，而远距离区域覆盖不足导致欠拟合。D$^2$GS 为此引入了一种新的评估指标来量化所学高斯分布的稳定性。

\subsection{扩散先验引导}
\label{ss:12_1b}

本节讨论两条技术路线：（1）扩散先验与伪视角生成；（2）扩散先验的变体与高效方法。前者关注如何利用扩散模型生成合理的伪视角来增强训练监督，后者关注如何降低扩散模型的计算开销。

\subsubsection{扩散先验与伪视角生成}

本节介绍利用扩散模型（Diffusion Model）生成伪视角来增强稀疏视角训练监督的一系列方法。扩散模型为稀疏视角 3DGS 提供了强大的视觉先验，可以生成合理的伪视角来增强训练监督。该路线的核心挑战在于如何确保生成内容的多视角一致性。

ReconX（2024.08）~\cite{2408_16767} 是该方向的早期代表性工作。

\textbf{基线与问题}：传统稀疏视角重建方法依赖深度、法线或代价体等几何约束来填充未见区域，但输入视角不足时这些 3D 先验难以在观测覆盖范围外生成高质量内容。

\textbf{核心改动}：ReconX 将稀疏视角重建重新定义为时序生成任务。它利用预训练视频扩散模型的强生成先验来创造额外观测，将欠约束的 3D 创建问题转化为充分约束的 3D 重建问题。

\textbf{实现原理}分为三阶段。（1）\textit{3D 结构条件构建}：利用 DUSt3R 从 $K$ 张稀疏输入图像（少至 2 张）恢复全局对齐点云 $P^*$，通过最远点采样（FPS）降至 $1/8$ 规模后，经 Transformer 编码器（列式正弦位置编码 + MLP + 交叉注意力）将其编码为 3D 上下文表示 $F(P)$，作为视频扩散的结构引导信号。该编码器与视频扩散模型端到端联合训练，确保其适配 DUSt3R 点云特性。（2）\textit{3D 一致性视频帧生成}：将首尾两个参考视角作为视频插值的首末帧，以图像条件 $F_g = g(I_{\mathrm{ref}})$（CLIP 编码）和 3D 结构条件 $F_F = F(P)$ 双重引导视频扩散模型。3D 条件通过 U-Net 空间层的交叉注意力注入：$F_{\mathrm{out}} = \mathrm{softmax}(Q^T K_g / \sqrt{d})V_g + \lambda_F \cdot \mathrm{softmax}(Q^T K_F / \sqrt{d})V_F$，其中 $\lambda_F$ 平衡图像条件与 3D 结构条件的贡献。训练目标采用无分类器引导（CFG）策略：$\mathcal{L}_{\mathrm{diffusion}} = \mathbb{E}_{x, p_\theta, \epsilon, t}[\| \epsilon - \epsilon_\theta(x_t, t, c_{\mathrm{view}}, c_{\mathrm{struc}}) \|_2^2]$。（3）\textit{置信度感知 3DGS 优化}：由于生成帧并非真实照片，直接用传统重建方法处理会产生伪影。ReconX 利用 DUSt3R 的逐像素置信度图 $C_i$ 衡量生成帧的可靠性，并引入 LPIPS 损失去除高层语义伪影，构建置信度加权损失：

\begin{equation*}
\mathcal{L}_{\mathrm{conf}} = \sum_{i=1}^{K'} C_i \bigl( \lambda_{\mathrm{rgb}} \mathcal{L}_1(\hat{I}^i, I^i) + \lambda_{\mathrm{ssim}} \mathcal{L}_{\mathrm{ssim}}(\hat{I}^i, I^i) + \lambda_{\mathrm{lpips}} \mathcal{L}_{\mathrm{lpips}}(\hat{I}^i, I^i) \bigr),
\end{equation*}

其中置信度图等效于不确定性度量，完成全局对齐的同时抑制不可靠像素的影响。

\textbf{设计动机}：从分布优化的视角，ReconX 证明了引入原生 3D 先验 $F$ 的条件扩散方案 $p_{\theta^*, \phi^*=\{g^*, F^*\}}(x)$ 比纯图像条件方案 $p_{\theta^*, \phi^*=g^*}(x)$ 更接近真实渲染分布 $q(x)$，即 $D(q(x) \| p_{\theta^*, \phi^*=\{g^*, F^*\}}(x)) < D(q(x) \| p_{\theta^*, \phi^*=g^*}(x))$，这为 3D 结构条件的引入提供了理论依据。

随后，IPSM-Gaussian（2024.12）~\cite{2412_02225} 深入分析了 Score Distillation Sampling（SDS，分数蒸馏采样，一种利用预训练扩散模型的梯度来指导 3D 参数优化的技术）在稀疏视角下的模式偏差（mode collapse，即优化结果被扩散模型的先验分布主导而丧失场景特异性）问题，提出内联先验引导的分数匹配（IPSM），利用视角间的姿态关系提供视觉内联先验来校正 SDS 的优化方向。

在扩散模型的高效利用方面，RealisticDreamer（2025.11）~\cite{2511_11213} 提出引导分数蒸馏（GSD，Guided Score Distillation），结合深度翘曲引导和语义特征引导来校正视频扩散模型的噪声预测方向，确保分数更新与正确的相机姿态和几何一致。FixingGS（2025.09）~\cite{2509_18759} 提出一种无需训练的方法，直接利用现有扩散模型为稀疏视角 3DGS 提供更准确的跨视角一致先验。GSFixer（2025.08）~\cite{2508_09667} 构建了基于 DiT（Diffusion Transformer）的视频扩散修复模型，以参考视角的 2D 语义特征和 3D 几何特征为条件，修复 3DGS 渲染的伪影帧。

VidSplat（2026.05）~\cite{2605_11424} 是一个训练无关的生成重建框架，利用视频扩散先验迭代合成新视角来补偿输入覆盖不足，并提出几何引导的交叉视角注意力机制（GCA，Geometry-guided Cross-Attention），通过预测深度图和相机姿态构建几何对齐的代理查询来替代损坏的渲染特征。GeoQuery（2026.05）~\cite{2605_12399} 同样采用几何引导的扩散框架，但通过几何诱导的对应场采样参考特征，并在局部窗口内限制交叉注意力以抑制伪匹配。OracleGS（2025.09）~\cite{2509_23258} 提出了一种新颖的"提出-验证"框架：先利用 3D 感知扩散模型合成新视角来提议完整场景，再利用多视角立体模型作为 3D 感知的预言机来验证生成视角的不确定性，以不确定性加权损失引导 3DGS 优化。

\subsubsection{扩散先验的变体与高效方法}

为了降低扩散模型引入的计算开销，研究者们探索了多种高效方案。本节按频率域解耦、潜空间操作和视频扩散先验三个角度组织这些方法。

WaveletGaussian（2025.09）~\cite{2509_19073} 系统性地将扩散过程从 RGB 域转移到小波（wavelet）域，其设计思路基于一个关键观察：扩散模型擅长修复低频结构信息但不善于处理高频纹理细节，因此将两者解耦处理更为合理。具体而言，该方法通过正向离散小波变换（DWT，Discrete Wavelet Transform）将 2D 图像分解为四个子带：LL（低频$\times$低频，保留粗结构）、LH（低频$\times$高频，捕获水平边缘）、HL（高频$\times$低频，捕获垂直边缘）和 HH（高频$\times$高频，强调对角纹理）。在训练流程上，WaveletGaussian 采用四阶段策略。（1）\textbf{粗训练}：在全部 $N$ 个稀疏视角上快速训练粗 3DGS 模型 $G_c$ 以获取大致几何，由于提前终止训练，即使已知视角的渲染也会出现中等程度的退化。（2）\textbf{数据集构建}：引入在线随机掩码（ORM）策略替代传统 leave-one-out 方法。

仅训练一个 $G_d$ 模型，对参考图像施加具有正弦位移的二值掩码以模拟多样化的损坏模式，从而避免为每个视角单独训练 3DGS 模型的巨大开销。将损坏-干净图像对变换到小波域，分别构造 LL 子带的扩散数据集和 LH/HL/HH 高频子带的 U-Net 修复数据集。（3）\textbf{扩散微调}：在低分辨率 LL 域中微调预训练扩散模型 $D$（本质上是低频修复的 inpainting 模型），同时微调轻量 U-Net $U$ 修复高频子带。两个模型各司其职且均工作在半分辨率下，计算成本远低于在 RGB 域微调单一扩散模型。（4）\textbf{精训练}：冻结的 $D$ 修复新视角渲染的 LL 子带，冻结的 $U$ 修复高频子带，两者通过逆 DWT 映射回原始分辨率的 RGB 域，作为伪参考提供无伪影的密集视角监督。通过频率解耦设计，WaveletGaussian 在保持渲染质量的同时大幅减少训练时间。

DWTGS（2025.07）~\cite{2507_15690} 同样利用小波空间损失进行频率正则化，但专注于低频子带的监督和高频子带的稀疏约束，以减少高频幻觉。LGDWT-GS（2026.01）~\cite{2601_17185} 则结合全局和局部频率正则化来稳定稀疏视角下的几何结构。

RI3D（2025.03）~\cite{2503_10860} 将视角合成过程分解为可见区域的"修复"（repair）和缺失区域的"补全"（inpainting）两个任务，分别训练两个个性化扩散模型，通过两阶段优化策略在可见和不可见区域均产生精细纹理。Splatent（2025.12）~\cite{2512_09923} 在 VAE（Variational Autoencoder，变分自编码器）潜空间中操作 3DGS，通过多视角注意力机制从输入视角的 2D 恢复细粒度细节，而非在 3D 空间中重建。

在视频扩散先验的利用方面，Taming Video Diffusion Prior（2025.03）~\cite{2503_05082} 解决了两个被忽视的实际问题：外推和遮挡。它提出了基于渲染序列的场景接地引导来驯服扩散模型生成一致的帧序列。FlowR（2025.04）~\cite{2504_01647} 提出多视角流匹配模型，直接学习从稀疏重建渲染到密集重建渲染的映射，在 360 万图像对上训练后可一次前向传播处理多达 45 个视角。BetterScene（2026.02）~\cite{2602_22596} 在 Stable Video Diffusion 框架中引入时间等变正则化和视觉基础模型对齐的表示，以增强 VAE 模块对 3D 一致性的建模。

ArtifactWorld（2026.04）~\cite{2604_12251} 建立了 3DGS 伪影的精细分类学，构建了包含 107.5K 多样化配对视频片段的训练集，通过同构预测器定位结构缺陷并引导修复。3DGS-Enhancer（2024.10）~\cite{2410_16266} 利用 2D 视频扩散先验解决 3D 视角一致性问题，通过时空解码器整合增强特征，并微调初始 3DGS 模型。

\subsection{初始化改进与综合优化策略}
\label{ss:12_2}

本节讨论两条技术路线：（1）初始化改进与致密化策略；（2）综合框架与特定场景方法。

\subsubsection{初始化改进与致密化策略}

本路线介绍改进稀疏视角下点云初始化质量与高斯致密化策略的方法。

初始化质量对稀疏视角 3DGS 的最终性能具有决定性影响。稀疏点云会导致高斯无法覆盖场景空缺，进而引发严重的过拟合。本节首先详细介绍里程碑式工作 FSGS，然后按稠密初始化、致密化策略两条线梳理后续进展。FSGS（2023.12）~\cite{2312_00451} 是该方向的里程碑式工作，其核心创新在于系统性地解决了标准 3DGS 在极稀疏点云初始化下的两大失败模式：（1）启发式位置梯度致密化无法覆盖整个场景空缺，以及（2）部分高斯趋向生长为极大体积导致过拟合训练视角。FSGS 提出三管齐下的解决方案：

（1）\textbf{邻近引导的高斯反池化（Proximity-guided Gaussian Unpooling）}：受计算机图形学中网格细分算法的启发，该方法在优化过程中动态构建邻近图（proximity graph），连接每个高斯与其 $K$ 个最近邻（实践中 $K{=}3$），并计算邻近分数 $P_i = \frac{1}{K}\sum_{j=1}^{K} D_i^K$，即高斯 $G_i$ 到其 $K$ 近邻的平均欧氏距离。当 $P_i$ 超过阈值 $t_p$ 时，在源高斯与目标高斯连线的中点生成新高斯，其尺度与不透明度属性继承目标高斯，旋转和球谐系数则初始化为零。这一策略鼓励新致密化的高斯分布在最具代表性的位置周围，在优化过程中逐步填充观测空缺，而非仅在已有高斯附近分裂或复制。

（2）\textbf{伪视角增强}：为缓解稀疏视角下视差不足导致的过拟合，FSGS 在训练视角间通过插值相机朝向和添加高斯位置噪声 $\mathcal{N}(0, \sigma)$ 来合成伪相机视角。这些在线采样的伪视角在 2000 次迭代后启用（确保高斯已大致表达场景结构），使 3D 高斯在优化过程中持续接收新视角的几何约束，从而降低对训练视角的过拟合风险。

（3）\textbf{单目深度引导的几何正则化}：FSGS 利用预训练的 Dense Prediction Transformer（DPT）估计训练视角和伪视角的单目深度图，并通过 Pearson 相关系数 $\rho(\mathbf{D}_{ras}, \mathbf{D}_{est}) = \frac{\mathrm{Cov}(\mathbf{D}_{ras}, \mathbf{D}_{est})}{\sqrt{\mathrm{Var}(\mathbf{D}_{ras}) \cdot \mathrm{Var}(\mathbf{D}_{est})}}$ 衡量渲染深度与估计深度的分布差异。这种放松的约束允许对齐深度结构而不受绝对深度值尺度不一致的阻碍。为实现端到端可微，FSGS 实现了可微深度光栅化器，利用与 3DGS 颜色渲染相同的 $\alpha$-混合公式累积有序高斯的 z-buffer 值。总训练损失为 $\mathcal{L} = 0.8 \mathcal{L}_1 + 0.2 \mathcal{L}_{ssim} + 0.05 \mathcal{L}_{regularization}$，其中几何正则化项同时作用于训练视角和伪视角。

FSGS 仅需 3 张训练视角即可实现实时、照片级的新视角合成，成为该方向的基准方法。

\textbf{局限。}FSGS 依赖预训练 DPT 提供几何正则化，Pearson 相关系数仅约束深度分布结构而缺乏对绝对尺度的显式建模，在深度尺度模糊场景中几何精度受限；伪视角增强固定 2000 次迭代后启用，难以适应不同场景的收敛速度。

在稠密初始化方面，Dust-GS（2024.09）~\cite{2409_08613} 引入自适应深度掩码技术来增强稀疏视角下的点云初始化质量。LoopSparseGS（2024.08）~\cite{2408_00254} 提出基于循环的渐进高斯初始化策略（PGI），利用训练过程中渲染的伪图像迭代致密化初始点云，同时通过深度对齐正则化（DAR）提供精确的几何监督。CoherentGS（2024.03）~\cite{2403_19495} 引入结构化高斯表示，通过隐式卷积解码器和全变差损失约束高斯位置，防止其在优化过程中独立移动。TWINGS（2026.05）~\cite{2605_22069} 利用薄板样条（Thin Plate Splines）进行非刚性变形建模，将反投影点与三角化控制点对齐以获得几何精确的初始化。

在致密化策略方面，AD-GS（2025.09）~\cite{2509_11003} 发现不受控的致密化是导致稀疏视角伪影的关键因素，提出交替致密化框架，在高致密化阶段（激进添加基元）和低致密化阶段（激进剪枝+几何正则化）之间交替。SDI-GS（2025.09）~\cite{2509_11853} 利用分割驱动初始化，仅保留结构显著区域的点来减少高斯数量。Opacity-Gradient Driven Density Control（2025.10）~\cite{2510_10257} 用不透明度梯度替代位置梯度作为致密化触发器，在 3 视角 LLFF 上比 FSGS 紧凑 40\% 以上。VGNC（2025.04）~\cite{2504_14548} 首次利用生成式新视角合成模型产生验证图像，据此确定最优高斯数量以减轻过拟合。

\subsubsection{综合框架与特定场景方法}

多种综合框架将上述正则化、扩散先验、初始化改进等技术路线整合为统一管道，同时也有工作针对特定场景（全景、不受约束照片集、无界场景等）设计了专门方案。COSMOS（2026.02）~\cite{2602_06044} 引入``超 Gaussian''（Super Gaussian）概念，基于局部几何和外观特征将高斯分组，在组间应用全局自注意力、组内应用稀疏局部注意力来集成空间先验。该方法受 3D 分割中超点（superpoint）概念的启发，利用结构感知特征预测高斯属性以促进一致的 3D 重建，同时通过组内位置正则化维持结构连贯性并抑制浮游物，在无需外部深度监督的条件下超越了当时最优方法。HeroGS（2026.03）~\cite{2603_01099} 建立图像级、特征级和参数级的层次化引导框架：图像级将稀疏监督转化为伪稠密引导，特征级利用低级特征细化高频细节，参数级通过共剪枝几何一致性去除不一致的 splat。DOC-GS（2026.04）~\cite{2604_06739} 从双域（优化域+观测域）视角建模高斯基元的可靠性，提出连续深度引导的 dropout 和基于暗通道先验的可靠性驱动几何剪枝。

在特定场景方面，PanoPlane（2026.05）~\cite{2605_14135} 针对室内 360 度全景场景，提出布局锚定注意力引导机制，将扩散模型的注意力导向检测到的平面表面，仅用 3 个输入视角即可实现高质量新视角合成。SparseGS-W（2025.03）~\cite{2503_19452} 和 MS-GS（2025.09）~\cite{2509_15548} 分别处理不受约束的真实世界照片集中的稀疏视角问题，前者利用约束扩散先验处理遮挡和外观变化，后者引入多外观能力以适应不同时间和季节的拍摄条件。Free360（2025.03）~\cite{2503_24382} 提出分层高斯表示来解决无界 360 度场景的极稀疏无位姿重建。Pano2Room（2024.08）~\cite{2408_11413} 从单张全景图重建室内 3D 场景，通过迭代全景 RGBD 修复逐步收集伪新视角来训练 3DGS。

NexusGS（2025.03）~\cite{2503_18794} 利用核线几何（epipolar geometry）将深度信息直接嵌入点云，通过核线深度连接、流鲁棒深度混合和流过滤深度剪枝三步构建稠密初始点云。GBR（2024.12）~\cite{2412_05908} 集成神经束调整模块和生成式深度精修模块，仅用 4--6 张输入视角即可实现高保真的高斯泼溅和网格重建。CuriGS（2025.11）~\cite{2511_16030} 引入课程学习框架，在训练视角周围采样不同扰动级别的学生视角，渐进解锁更高扰动级别以稳定增强训练。S2Gaussian（2025.03）~\cite{2503_04314} 同时处理稀疏视角和低分辨率的复合退化问题，通过两阶段策略先优化低分辨率高斯表示再精修高分辨率版本。

此外，若干工作从信息论角度分析稀疏视角问题。UGOD（2025.08）~\cite{2508_04968} 引入学习式不确定性（learnable uncertainty，即通过网络自身预测每个高斯重建可靠性的置信度分数）来引导可微不透明度更新和软 dropout 正则化。Co-Adaptation 分析（2025.08）~\cite{2508_12720} 揭示了稀疏视角 3DGS 中的共适应效应（高斯基元之间过度纠缠以拟合训练视角），提出 Co-Adaptation Score 度量并提出随机 dropout 和乘性噪声注入两种轻量缓解策略。CoMapGS（2025.03）~\cite{2503_20998} 利用共可见性图（covisibility map）来识别和增强欠表示的稀疏区域。GSCompleter（2026.04）~\cite{2604_20155} 提出``生成-注册''工作流替代不稳定的``修复-蒸馏''范式，通过立体锚点视角选择和射线约束注册将新视角图像提升为 3D 高斯基元。

其他重要贡献包括：SCGaussian（2024.11）~\cite{2411_03637} 引入混合高斯表示（非结构化 + 射线绑定，即部分高斯独立优化、部分高斯约束在相机射线上），利用匹配对应关系直接约束高斯位置收敛到表面点。MCGS（2024.10）~\cite{2410_11394} 利用稀疏匹配器的匹配先验来初始化高斯，并提出多视角一致性引导的渐进剪枝策略。Binocular-GS（2024.10）~\cite{2410_18822} 无需外部先验，通过视差引导的图像翘曲构造双目立体对，利用自监督的立体一致性来约束 3DGS 优化。EG-Gaussian（2025.04）~\cite{2504_13540} 将核线几何融入 3DGS 初始化阶段，并设计图学习模块精修空间特征。

FewViewGS（2024.11）~\cite{2411_02229} 通过多阶段训练方案和基于匹配的一致性约束来增强新视角渲染。SPARS3R（2024.11）~\cite{2411_12592} 结合 SfM 的精确位姿估计和深度估计的稠密点云，通过语义离群值对齐进一步提升稀疏视角重建质量。PointGS（2025.06）~\cite{2506_10335} 利用立体基础模型估计位姿和稠密点云，并通过点交互网络增强逐点外观表示。NovelGS（2024.11）~\cite{2411_16779} 利用扩散模型进行新视角去噪来生成 3D 高斯。

SVS-GS（2024.09）~\cite{2409_03213} 集成 3D 高斯平滑滤波器和深度梯度轮廓先验损失，利用几何平滑性约束来抑制稀疏视角下的浮游物和几何退化。LM-Gaussian（2024.09）~\cite{2409_03456} 利用大型视觉模型先验来增强稀疏视角重建。MVPGS（2024.09）~\cite{2409_14316} 利用学习式多视角立体的先验来增强几何初始化和外观约束。AugGS（2024.08）~\cite{2408_04831} 通过自增强两阶段框架和结构掩码策略来优化稀疏视角重建。FSFSplatter（2025.10）~\cite{2510_02691} 实现端到端的稠密高斯初始化和几何增强场景优化。Sparse2DGS（2025.04）~\cite{2504_20378} 提出几何优先的增强方案用于稀疏视角表面重建，将二维高斯 surfel（surface element，表面元素）引入稀疏视角设定。BRUM（2025.07）~\cite{2507_12095} 针对车辆的 360 度稀疏图像重建进行了专门优化。Sparse2DGS（2025.05）~\cite{2505_19854} 利用 DUSt3R 和 COLMAP MVS 生成精确稠密的 3D 点云。Intern-GS（2025.05）~\cite{2505_20729} 利用视觉基础模型引导初始化和优化过程。

PointmapDiff（2025.01）~\cite{2501_02913} 利用点图作为扩散模型的条件信号。TSGaussian（2024.12）~\cite{2412_10051} 结合语义约束和深度先验来重建特定目标。SPC-GS（2025.03）~\cite{2503_12535} 利用场景布局初始化和语义提示一致性约束。HBSplat（2025.09）~\cite{2509_24893} 提出混合损失深度估计和双向翘曲虚拟视角合成。FrameTwin（2026.05）~\cite{2605_09362} 将高斯锚定到参数化曲线上进行线框结构的稀疏视角捕获。S2D（2026.03）~\cite{2603_10893} 通过稀疏到稠密的提升管道，利用一步扩散模型修复图像伪影。Sparse-View 3DGS in the Wild（2026.04）~\cite{2604_27422} 处理含有干扰物的真实世界稀疏视角场景。Gaussian Scenes（2024.11）~\cite{2411_15966} 利用深度增强的扩散先验进行无位姿稀疏视角重建。Init-to-Generalize（2025.10）~\cite{2510_17479} 证明初始化是稀疏视角 3DGS 的决定性因素。

Seeing 3D World（2025.03）~\cite{2503_00260} 设计了折反射成像系统从单次拍摄中恢复微型场景的 3D 视图。Few-shot 3DGS survey（2025.07）~\cite{2507_16406} 对该领域进行了系统性综述。GaussianCAD（2025.03）~\cite{2503_05161} 将 CAD 重建视为稀疏视角 3D 重建的特例。Dust to Tower（2024.12）~\cite{2412_19518} 提出由粗到精的框架从稀疏未标定图像同时优化 3DGS 和位姿。

\begin{takeaway}
稀疏视角新视角合成的核心挑战是严重的过拟合问题，现有方法沿正则化约束、扩散模型伪视角增强、初始化改进三条路线展开，其中几何基础模型（DUSt3R 等）正逐步取代传统 SfM 成为稀疏视角重建的核心基础设施。
\end{takeaway}

% ======================================================================

\subsection{近期方法脉络：稀疏视角几何约束、不确定性与信息量指导、Sparse panoramic reconstruction等}

这一组近期工作延续 Gaussian Splatting、3DGS，主要把改动集中在正则化、优化、增量改进等环节。下面只保留相对既有路线的变化，便于和前文方法直接对照。

\paragraph{VisDom}
VisDom（2026.06）~\cite{2606_20531} 延续 Gaussian Splatting 在“稀疏视角几何约束”上的基本设定，主要改变是：VisDom：Visible Domain Constraint，基于visual hull carving加入最小多视角可见域约束。3DGS 的显式高斯表示和可微渲染能力保持为基础。
其设计动机是：提出无需学习的几何约束，弥补仅用silhouette一致性会扩展到真实几何外部的问题。
具体而言，方法在经典visual hull基础上强制最小多视角可见域，作为NeRF或Gaussian Splatting稀疏重建的正则。
适用边界是：依赖轮廓/可见域估计；对复杂拓扑和遮挡可能保守。

\paragraph{From Uncertainty to Stability and Fidelity}
From Uncertainty to Stability and Fidelity（2026.06）~\cite{2606_20842} 延续 3DGS 在“不确定性与信息量指导”上的基本设定，主要改变是：基于Fisher Information分析视角/参数不确定性，引导novel-view采样和Dropout式正则。3DGS 的显式高斯表示和可微渲染能力保持为基础。
其设计动机是：将Fisher信息引入稀疏3DGS，替代随机采样或经验式Dropout，为不确定区域提供可解释训练信号。
具体而言，方法估计3DGS参数或视角对重建的Fisher信息，识别高不确定区域并指导新视角增强、正则化或采样策略。
适用边界是：Fisher估计增加计算；近似质量影响指导效果。

\paragraph{PanoImager}
PanoImager（2026.06）~\cite{2606_27071} 延续 3DGS 在“Sparse panoramic reconstruction”上的基本设定，主要改变是：前馈位姿深度先验、几何条件扩散补视图和深度引导 3DGS 优化。3DGS 的显式高斯表示和可微渲染能力保持为基础。
其设计动机是：提出无需 SfM 的稀疏全景重建框架，解决旋转主导运动下初始化不可靠的问题。
具体而言，将全景图拆成局部透视视图，使用 feed-forward pose/depth prior 和 diffusion view completion 生成辅助观测，再通过 depth-guided 3DGS 优化提升跨视角一致性。
适用边界是：依赖深度/位姿先验和生成式补视图，可能引入 hallucination；全景相机模型和遮挡处理是关键风险。

\paragraph{StereoGS}
StereoGS（2026.06）~\cite{2606_30545} 延续 3DGS 在“Stereo priors for sparse-view Gaussian Splatting”上的基本设定，主要改变是：构造虚拟 stereo pairs 并使用 stereo depth regularization 建立双目一致性。3DGS 的显式高斯表示和可微渲染能力保持为基础。
其设计动机是：在优化中构造虚拟双目对，引入更可靠的 binocular consistency，提高稀疏输入下的几何约束。
具体而言，相比 monocular depth priors，StereoGS 的 stereo depth regularization 提供跨视图尺度一致的深度监督，减少 sparse-view 3DGS 的 defective geometry。
适用边界是：立体先验质量和虚拟视角构造会影响效果；极少视角和弱纹理区域仍可能不稳。

\paragraph{ACEsplat}
ACEsplat（2026.06）~\cite{2606_22091} 延续 3DGS 在“无外部几何初始化”上的基本设定，主要改变是：ACEsplat：自监督scene coordinate regression构建内部几何先验，再进行加速3D Gaussian scene regression。3DGS 的显式高斯表示和可微渲染能力保持为基础。
其设计动机是：提出无需预计算SfM或深度的两阶段快速3DGS优化流程，提升现场捕获部署可靠性。
具体而言，第一阶段SCR模块从RGB和位姿自监督学习场景坐标并建立内部几何先验；第二阶段用该先验加速高斯场优化。
适用边界是：仍需已知相机位姿；SCR先验失败会影响后续高斯优化。

\paragraph{AugSplat}
AugSplat（2026.06）~\cite{2606_31556} 延续 3DGS 在“Radiance-field priors for sparse-view Gaussian Splatting”上的基本设定，主要改变是：Radiance Field-informed Gaussian Splatting，用 NeRF/RF 信息增强稀疏视图 3DGS 初始化和优化。3DGS 的显式高斯表示和可微渲染能力保持为基础。
其设计动机是：结合 NeRF 类方法在少视角鲁棒重建上的优势和 Gaussian Splatting 的实时渲染优势。
具体而言，在稀疏输入下，利用 radiance field 提供更稳健的几何或外观线索，缓解不完整、噪声点云初始化导致的 Gaussian 优化失败。
适用边界是：可能引入额外 radiance field 训练成本；稀疏极限和复杂遮挡场景仍需要验证。

\paragraph{Rectifying Mask via Entropy for Distractor-Free 3DGS in Ambiguous Scenarios}
Rectifying Mask via Entropy for Distractor-Free 3DGS in Ambiguous Scenarios（2026.06）~\cite{2606_29496} 延续 3DGS 在“Distractor-robust Gaussian reconstruction”上的基本设定，主要改变是：基于熵的自适应 transient mask 与 density control。3DGS 的显式高斯表示和可微渲染能力保持为基础。
其设计动机是：针对颜色或语义歧义下难以区分 transient 元素和静态场景的问题，提出系统化 mask 构建和密度控制。
具体而言，结合熵信号和实例 mask 捕获模糊干扰物，并用 entropy-aware positional gradients 调整 Gaussian density control，使高斯更好对齐真实静态结构。
适用边界是：依赖熵估计和实例 mask 质量，对长期静止的动态物体或语义错误仍可能失效。

\paragraph{Improving Sparse-View 3DGS Generalization via Flat Minima Optimization}
Improving Sparse-View 3DGS Generalization via Flat Minima Optimization（2026.07）~\cite{2607_00885} 延续 3DGS 在“Flat-minima optimization for sparse-view 3DGS”上的基本设定，主要改变是：将 flat minima optimization 原理适配到 Gaussian 参数训练。3DGS 的显式高斯表示和可微渲染能力保持为基础。
其设计动机是：把稀疏视图过拟合解释为尖锐极小值问题，并提出轻量训练框架寻找更平坦的解。
具体而言，将 Gaussian positions、covariances、opacity 和 appearance 等参数视作可训练权重，在优化中加入 flat minima 风格的扰动稳定性约束。
适用边界是：平坦极小值假设在几何参数上的适用边界需要更多理论分析；扰动尺度选择可能敏感。
\section{位姿估计与几何先验引导}
\label{sec:12_pose}

3DGS 的标准流程依赖 COLMAP 等 SfM 系统提供相机位姿和初始稀疏点云。然而，在稀疏视角、纹理缺失、大基线运动等条件下，SfM 往往失败或产生不准确的位姿，直接导致 3DGS 训练崩溃。此外，单目深度估计模型提供了丰富的几何先验，可在稀疏视角条件下为 3DGS 提供额外的几何监督。本节从位姿估计和几何先验引导两个角度展开讨论，涵盖联合位姿-外观优化、前馈式位姿估计、视频连续位姿追踪、特定场景的无 SfM 方案，以及深度正则化与对齐和几何约束与表面重建等子方向。

\subsection{位姿估计}
\label{ss:12_3}

本节讨论四条技术路线：（1）联合位姿-外观优化；（2）前馈式位姿估计；（3）视频连续位姿追踪；（4）特定场景与模态。

\subsubsection{联合位姿-外观优化}

本路线介绍在 3DGS 训练过程中同时优化高斯参数和相机位姿的联合优化方法。

联合优化的核心思想是在 3DGS 训练过程中同时优化高斯参数和相机位姿，通过可微渲染的反向传播来精修初始位姿。直觉上来说，如果渲染图像与真实图像存在差异，那么这个差异不仅可能源于高斯参数不够准确，也可能源于相机位姿估计有偏差。因此同时优化两者才能更有效地消除误差。

针对 3DGS 原版管线仅支持固定相机位姿下高斯参数优化、光栅化器缺失反向传播至 SE(3) 外参通道这一根本缺陷，早期工作"Look Gauss, No Pose"（2024.10）~\cite{2410_08743} 系统性地解决了 3DGS 光栅化器不提供相机外参梯度这一根本缺陷。具体而言，该方法将相机位姿建模为 SE(3) 李群元素（SE(3) 是描述三维刚体运动的数学群，包含三维旋转和三维平移共六个自由度），在对应的李代数 $\mathfrak{se}(3) \simeq \mathbb{R}^6$ 切空间中推导解析梯度。通过链式法则，光度损失对相机位姿的梯度由三个分量贡献。2D 投影均值通过刚体变换 Jacobian 传递梯度；2D 投影协方差通过旋转变换 Jacobian 贡献梯度；球谐函数颜色项则通过相机中心到高斯中心的方向向量传递梯度。这些梯度计算被完整集成到 CUDA 光栅化器的反向传播通道中，仅引入极小的额外计算开销。

在联合优化策略上，该方法采用 $L_1$ + DSSIM 的组合光度损失，并引入各向异性正则化（限制高斯长短轴比率，防止高斯退化为薄片状结构以"欺骗"渲染损失）。同时设计自适应透明度阈值剪枝策略，通过在训练前 10k 步添加 opacity $L_1$ 损失来减少半透明高斯的数量，缓解形状-辐射度模糊问题。此外，该方法还展示了从视频流中估计相机轨迹的能力：对相邻帧用单目深度估计器快速构建帧级高斯表示（不到 100 步优化，亚秒级完成），再用置信度掩码的光度损失精修帧间相对位姿，仅需部分重叠的几何假设即可完成全局轨迹估计。

GloSplat（2026.03）~\cite{2603_04847} 提出了联合位姿-外观优化框架，其关键创新在于将显式 SfM 特征轨迹（feature tracks，即同一三维点在多视角中的二维观测序列）作为一等实体在整个训练过程中维护。轨迹 3D 点作为独立于高斯基元的可优化参数，通过重投影损失与光度监督并行运作，防止早期位姿漂移。GloSplat 提供了两个变体：无需 COLMAP 的 GloSplat-F（基于检索的配对选择）和追求最高质量的 GloSplat-A（穷举匹配）。与依赖纯光度梯度的位姿精修方法不同，GloSplat 的特征轨迹提供了持久的几何锚点。实验表明 GloSplat-F 在无需 COLMAP 方法中达到当时最优，而 GloSplat-A 则超越了所有基于 COLMAP 的基线方法。3R-GS（2025.04）~\cite{2504_04294} 基于 MASt3R-SfM 的大型重建先验，通过优化的实践方案解决了朴素联合优化的两个挑战：对 SfM 初始化质量的敏感性和全局优化能力不足。

JOGS（2025.10）~\cite{2510_26117} 提出将联合优化解耦为两个交替阶段：固定位姿优化高斯参数，以及利用定制的 3D 光流算法结合几何和光度约束精修位姿。PCR-GS（2025.07）~\cite{2507_13891} 针对复杂相机轨迹场景，从特征重投影正则化（利用 DINO 特征的语义对齐）和小波频率正则化（利用高频细节差异优化旋转矩阵）两个角度进行位姿共正则化。

\subsubsection{前馈式位姿估计}

本路线介绍通过神经网络直接从图像预测相机位姿的前馈式方法。

前馈方法通过神经网络直接从图像预测相机位姿，避免了迭代优化的计算开销。与联合优化方法不同，前馈方法在单次前向传播中即可输出位姿估计，因此在推理速度上具有显著优势。

针对 SfM 依赖型 3DGS 在稀疏视角下特征匹配稀缺、位姿估计不可靠的局限，InstantSplat（2024.03）~\cite{2403_20309} 是该方向的里程碑工作，提出了一个完全自监督的端到端优化框架，将整个流程从 COLMAP（一种基于特征匹配的运动恢复结构系统）的分钟级缩短至秒级，实现超过 30 倍的加速（SSIM 从 0.3755 提升至 0.7624）。其核心管道包含三个关键组件。首先是共可见全局几何初始化：利用 MASt3R 几何基础模型对输入图像进行成对推理，回归像素对齐的点图、相机内参和外参；随后构建完全连通图将成对预测对齐到全局坐标系，并通过焦距平均化稳定单相机假设下的内参估计。其次是共可见冗余消除：按视角平均置信度排序，从低置信度视角出发，将高置信度视角的点图投影过来进行深度验证，差值低于阈值的点被判定为冗余并剪枝——这一步骤显著减少了自由参数数量并加速了后续优化。第三是置信度感知的高斯束调整：在端到端优化中同时精修高斯基元参数 $G$ 和相机位姿 $T$，核心损失函数为纯光度误差（$L_1$ + SSIM），不依赖任何启发式正则化；其置信度感知优化器根据 MASt3R 输出的逐点置信度动态调整学习率——对低置信度（初始几何误差大）的点赋予更高学习率，确保优化资源集中到最需要修正的区域。整个流程从无位姿的原始 RGB 图像出发，通过前馈推理获得初始几何和位姿，再经少量步数的梯度优化即可达到高质量重建。

iGaussian（2025.11）~\cite{2511_14149} 提出首个实时前馈位姿估计框架，通过高斯场景先验的位姿回归网络预测粗 6-DoF 位姿，再通过特征匹配和多模型融合进行精修，实现了 2.87 FPS 的追踪速度和仅 0.2 度的中位旋转误差。

SPFSplat（2025.08）~\cite{2508_01171} 及其升级版 SPFSplatV2（2025.09）~\cite{2509_17246} 提出了自监督无位姿框架：共享特征提取骨干同时预测 3D 高斯基元和典型空间（canonical space，一个统一的三维参考坐标系）中的相机位姿，通过重投影损失约束像素对齐的高斯基元。Gesplat（2025.10）~\cite{2510_10097} 利用 VGGT 基础模型获得更可靠的初始位姿和稠密点云，并提出混合高斯表示、图引导属性精修和基于流的深度正则化。

RegGS（2025.07）~\cite{2507_08136} 提出 3D 高斯注册（registration，将两个三维表示对齐到统一坐标系的操作）框架，利用熵正则化 Sinkhorn 算法求解高斯混合模型的最优传输距离作为对齐度量，实现从粗到精的注册过程。Probabilistic Procrustes Mapping（2025.07）~\cite{2507_18541} 将大规模户外场景的位姿估计转化为概率 Procrustes 问题，通过软 dustbin 机制拒绝不确定的对应关系，在数分钟内对齐数百张图像的点云和位姿。

\subsubsection{视频连续位姿追踪}

本路线介绍针对视频输入处理位姿漂移和长序列连续追踪的方法。

针对视频输入，位姿漂移和长序列处理是核心挑战。视频的时序连续性为位姿估计提供了额外的约束——相邻帧之间的视角变化通常很小，可以被有效利用来稳定追踪。

Rob-GS（2025.01）~\cite{2501_15096} 提出鲁棒的渐进式框架，其核心设计包含三个紧密耦合的模块。首先，\textbf{流引导的邻接位姿追踪}利用视频帧的固有连续性——相邻帧之间的视角变化通常很小，可以被有效利用来稳定追踪。对每一帧先用单目深度估计反投影为彩色点云并拟合单帧专属高斯 $G_i$（不做剪枝或致密化），然后以 $G_i$ 渲染下一帧候选视角下的图像，联合最小化光度损失 $L_c$ 和流损失 $L_f$ 来优化位姿 $T_{i+1}$。其中流损失是该追踪策略的关键：通过 RAFT 计算帧间光流 $f_{i+1}$，同时根据当前深度图和相邻位姿计算重投影流 $\hat{f}_{i+1}$（将像素反投影至 3D 再投影到下一帧坐标系），以 $\|\hat{f}_{i+1} - f_{i+1}\|^2$ 作为监督信号，将稠密对应关系转化为对位姿的梯度约束，使追踪在大幅视角变化下仍保持鲁棒。此外，作者推导了 $\alpha$-混合 过程相对于 SE(3) 位姿参数的解析雅可比矩阵，修改了 CUDA 光栅化器以支持位姿梯度的反向传播。

其次，\textbf{可见性保持检查}策略以极低计算开销自适应分割视频序列：维护前一段末端帧的拟合高斯集合 $G_a$（含 $V_a$ 个高斯），对每个新追踪帧将其投影到该帧视角下统计可见高斯数 $V_i$，计算保持率 $R_c = V_i / V_a$。当相机缓慢推进时 $R_c$ 保持高位、不触发分割；当相机大幅平移或旋转至新区域时 $R_c$ 降至阈值以下，自动触发分割更新和局部优化。这一机制精准区分了渐进式视角变化与场景跳变。最后，每个局部段以首帧（必要时补充末帧）的单目深度反投影初始化 3DGS，在段内随机采样帧联合优化颜色损失和尺度不变深度损失，相邻段之间保持 5 帧重叠以确保无缝衔接。

LongSplat（2025.08）~\cite{2508_14041} 针对随意拍摄的长视频，引入增量联合优化（同时优化位姿和高斯）、鲁棒位姿估计模块和八叉树（octree，一种将三维空间递归划分为八个卦限的层次化数据结构）锚点形成机制。KeyGS（2024.12）~\cite{2412_20767} 以关键帧为中心，先用 SfM 快速获取粗位姿（秒级），再利用 3DGS 的稠密表示精修位姿，将训练时间从小时级压缩到分钟级。

SFGS（2024.12）~\cite{2412_01553} 引入层次化训练策略，将视频分为多个片段分别训练高斯表示，再合并为统一模型，利用视频帧插值模型补偿大相机运动。TrackGS（2025.02）~\cite{2502_19800} 引入全局特征轨迹建立全局几何约束，通过轨迹约束高斯、2D/3D 轨迹损失和可微相机内参优化实现无 COLMAP 的新视角合成。ZeroGS（2024.11）~\cite{2411_15779} 从数百张无序未位姿图像训练 3DGS，利用预训练基础模型作为神经场景表示，从种子图像逐步注册新图像。

FlowMap（2024.04）~\cite{2404_15259} 提出端到端可微方法，通过梯度下降最小化光流（optical flow，描述相邻帧之间像素运动的二维向量场）诱导的深度-内参-位姿一致性目标函数，令人惊讶地在 360 度新视角合成任务上达到了与 COLMAP 相当的效果。具体而言，FlowMap 利用现成的光流和点追踪获取帧间对应关系，并对深度、内参和位姿引入可微重参数化以适配一阶优化，同时利用点轨迹鼓励长期几何一致性，是一种完全脱离传统 SfM 管线的纯梯度下降方案。ICP-3DGS（2025.06）~\cite{2506_21629} 将迭代最近点（ICP）与优化精修结合，在大尺度室内外场景中实现准确的位姿估计。InstantSfM（2025.10）~\cite{2510_13310} 构建了全 GPU 原生的全局 SfM 系统，将度量深度先验嵌入全局定位和束调整中，在大规模场景上实现约 40 倍于 COLMAP 的加速。

\subsubsection{特定场景与模态}

本路线介绍事件相机、手术场景、水下环境等特定场景与模态下的无 SfM 方案。

在特定应用场景下，研究者提出了针对性的无 SfM（Structure from Motion，运动恢复结构）方案，涵盖事件相机、手术场景、水下环境、大视场角广角镜头等多种特殊设定。

首先，\textbf{基于块的时间相干性分析}从噪声事件流中稳健地提取边缘：事件相机（event camera，一种仅输出像素亮度变化事件的异步传感器，与传统相机输出完整帧图像根本不同）产生的原始事件流充满噪声，需要通过统计方法提取有意义的结构信息。将连续事件图划分为带重叠的空间块，对每个块位置提取跨时间步的块序列，计算经高斯窗处理的相邻事件图之间的时间差分信号 $D_t$。由于边缘在相机运动时在空间相干位置触发事件，产生高方差的结构化模式，而非边缘区域产生低方差的随机散布事件，因此通过最大化跨时间步的块对比度并配合自适应阈值即可稳健地分类边缘块。该设计还具备自纠错能力：时间差分同时捕获前后两帧的边缘信息，先前帧中深度估计错误的边缘可基于与当前观测的不一致性被识别和剔除。

其次，\textbf{边缘引导的高斯初始化}将检测到的边缘转化为结构化的 3D 高斯分布：从边缘图中提取 2D 边缘点后，通过 PCA（主成分分析，Principal Component Analysis，一种通过正交变换将数据投影到主成分方向的降维方法）确定边缘切线方向和法线方向，再使用递归网格策略生成对齐于边缘方向的 2D 边缘高斯。这些 2D 高斯通过\textbf{逆对数深度采样}扩展至 3D——沿观测射线以 $d = 1/(1/d_{\min} + u \cdot (1/d_{\max} - 1/d_{\min}))$（$u \sim \mathcal{U}(0,1)$）采样深度值，在远距离处密集采样。这一设计具有明确的几何动机：远距离点在相机旋转下产生更大的像素位移，因此对旋转运动的可观测性更强，密集采样远距离点显著提升了位姿估计的精度。剩余比例的高斯则随机初始化以覆盖纹理缺失区域。

最后，\textbf{边缘加权损失}在重建过程中根据边缘置信度空间加权光度误差：$w(x) = 1 + \beta \cdot M(x)$，其中 $M(x)$ 为边缘掩码、$\beta$ 为权重系数，使优化优先聚焦于几何显著性边界，实现更快的收敛和更准确的深度估计。EF-3DGS（2024.10）~\cite{2410_15392} 首次将事件相机引入随意拍摄视频的场景构建，通过事件生成模型（将事件流与亮度变化关联）、对比度最大化框架（从事件流中提取边缘信息）和光度束调整（联合优化相机位姿和场景结构）三个组件融合事件和帧信息。

AttentionGS（2025.06）~\cite{2506_23611} 消除了对高质量初始点云的依赖，从随机初始化出发，在早期训练阶段引入几何注意力（geometry attention，一种利用三维空间邻域关系约束高斯分布的注意力机制）快速恢复全局结构，后期引入纹理注意力精修细节。CSS（2024.09）~\cite{2409_08562} 针对众包图像，通过鲁棒几何先验和高级光照建模克服缺失位姿、有限视角和不一致照明的挑战。Correspondence-Guided SfM-Free（2024.08）~\cite{2408_08723} 利用目标与渲染结果之间的对应关系实现更好的像素对齐，避免不准确初始位姿导致的梯度不稳定问题。

Free-SurGS（2024.07）~\cite{2407_02918} 首次将无 SfM 的 3DGS 应用于手术场景重建，利用光流先验引导投影流并引入对极几何（epipolar geometry，描述两个视角之间点对应关系的几何约束）一致性检查过滤异常值。AquaGS（2025.05）~\cite{2505_01799} 专为水下场景设计，利用 SeaThru 算法快速分离场景细节和介质（水中悬浮颗粒、有色溶解有机物等导致散射和颜色偏移的因素）特征。Self-Calibrating GS（2025.02）~\cite{2502_09563} 提出自标定框架联合优化相机参数、镜头畸变和 3D 高斯表示，支持大视场角广角镜头的高质量重建。SmallGS（2025.04）~\cite{2504_17810} 专为小基线视频设计，将 DINOv2 高维特征融入 3DGS 以增强小基线下的位姿估计鲁棒性。TLC-Calib（2025.04）~\cite{2504_04597} 解决 LiDAR-相机无标定板的在线标定问题，冻结可靠 LiDAR 点作为锚点高斯以保持全局结构。

此外，GeGS-PCR（2026.04）~\cite{2604_17721} 利用颜色信息增强点云注册，提出几何-3DGS 模块编码局部邻域信息。GSA（2026.03）~\cite{2603_21936} 首次实现跨实例（同类别不同物体）的 3DGS 注册。PFGS（2025.10）~\cite{2510_15386} 处理多姿态物体采集的完整重建问题。Fast Learning of Spacecraft（2025.07）~\cite{2507_19459} 利用基于 CNN 的基元初始化器实现非合作航天器的快速 3D 建模。Adam SLAM（2025.08）~\cite{2508_20526} 通过 3DGS 反向传播精细化相机标定参数。Camera Pose Refinement via 3DGS（2025.08）~\cite{2508_17876} 利用多视角渲染的对极几何约束进行位姿精修。Rethinking Pose Refinement（2026.03）~\cite{2603_16538} 引入 Monte Carlo 位姿采样和 Fisher 信息引导的 PnP 优化来处理位姿和几何不确定性。Sparfels（2025.05）~\cite{2505_02178} 利用单一 3D 基础模型实现快速稀疏无位姿重建。Recollection from Pensieve（2025.05）~\cite{2505_13440} 提出两阶段策略从无标定视频训练新视角合成模型。Constrained Optimization（2025.04）~\cite{2504_09129} 结合 LiDAR 点云和粗位姿图像进行约束优化。SparseLGS（2024.12）~\cite{2412_02245} 首次在无位姿稀疏视角下解决 3D 语义场问题。Evaluating Alternatives to SfM（2024.04）~\cite{2404_12547} 系统评估了替代 SfM 初始化的各种策略。

此外，单目深度估计（Monocular Depth Estimation, MDE，单目深度估计）模型提供了丰富的几何先验，可以在稀疏视角条件下为 3DGS 提供额外的几何监督。然而，MDE 固有的尺度模糊性、多视角不一致性和局部几何不准确性，使得直接应用可能反而降低渲染性能。

\subsection{几何先验与约束}
\label{ss:12_4}

本节讨论两条技术路线：（1）深度正则化与对齐；（2）几何约束与表面重建。

\subsubsection{深度正则化与对齐}

本路线介绍利用单目深度估计先验对 3DGS 施加深度正则化与对齐约束的方法。

直接利用 MDE 预测作为伪真值进行 L2 深度损失监督是最朴素的方案，但由于尺度模糊性和局部误差，效果往往有限。直觉上来说，MDE 模型可以告诉我们场景中"哪些物体更近、哪些更远"的大致排序，但具体的深度数值可能存在系统性偏差，因此需要设计能够容忍这些偏差的正则化策略。本节首先详细介绍 UGOT 的最优传输方法，然后梳理其他深度对齐与正则化方案。

UGOT（2024.05）~\cite{2405_19657} 认识到 MDE（Monocular Depth Estimation，单目深度估计）预测在特定区域存在显著不确定性，引入不确定性估计并集成基于块的最优传输策略来补充传统 L2 深度损失。

\textbf{基线与问题}：传统稀疏视角 3DGS 使用逐像素 $L_2$ 深度损失将渲染深度与 MDE 预测对齐，但 MDE 模型在纹理缺失、遮挡边界和重复结构等区域存在不可忽略的预测误差。直接逐像素强监督会导致高斯基元向错误深度偏移，反而降低渲染质量。此外，由于 3DGS 的体积渲染过程中每个像素对应大量高斯基元，其中绝大部分权重集中于少数高斯，其余小权重高斯的深度容易在优化中漂移，进一步恶化深度估计的可靠性。

\textbf{核心改动}：UGOT 从两个维度改进深度监督。（1）引入基于扩散模型的 MDE 不确定性量化机制，降低高不确定区域深度预测的监督权重。（2）将深度对齐从逐像素 $L_2$ 损失提升为基于块（patch-wise）的最优传输分布匹配，在局部空间上下文中建模渲染深度分布与 MDE 预测的一致性，而非强制逐像素精确匹配。

\textbf{实现原理}包含四个关键组件。

（1）\textit{深度采样}：在 3DGS 的前向渲染中，每个像素对应 $N$ 个高斯基元。UGOT 按权重提取每个像素的前 $K$ 个高斯的深度 $d_{\mathrm{top}} = \{d_1, d_2, \ldots, d_K\}$ 和权重 $w_{\mathrm{top}} = \{w_1, w_2, \ldots, w_K\}$，然后使用 Gumbel-Softmax 采样从中采样 $n$ 个深度值：

\begin{equation*}
p_j = \frac{\exp((\log(w_j) + g_j) / \tau)}{\sum_{k=1}^{K} \exp((\log(w_k) + g_k) / \tau)}, \quad d^*_i = \sum_{j=1}^{K} p_j \, d_j, \quad i = 1, 2, \ldots, n,
\end{equation*}

其中 $g_j$ 为独立同分布的 Gumbel 噪声变量，$\tau$ 为温度参数控制分布软度。采样策略既减少了计算和存储开销，又缓解了小权重高斯深度漂移的问题。

（2）\textit{基于块的最优传输深度损失}：将像素级深度采样提升为 patch 级分布匹配。每次迭代随机选择大小为 $s \times s$（$s \in [s_1, s_2]$）的 patch，对 patch 内每个像素的采样深度取均值 $\bar{D}_{\mathrm{patch}} = \{\bar{d}_1, \bar{d}_2, \ldots, \bar{d}_n\}$，与对应 MDE 深度真值均值 $\bar{D}_{\mathrm{gt}}$ 之间的最优传输距离定义为：

\begin{equation*}
\mathrm{OT}_\epsilon(\bar{D}_{\mathrm{patch}}, \bar{D}_{\mathrm{gt}}) = \min_{T_i \in \Pi(\bar{D}_{\mathrm{patch}}, \bar{D}_{\mathrm{gt}})} \sum_i C_i T_i - \epsilon \sum_i T_i \log T_i,
\end{equation*}

其中 $C_i$ 为传输代价，$\epsilon$ 为熵正则化系数，$\Pi$ 为满足边际约束的可行传输方案集合。基于块的聚合有效增强了深度优化的鲁棒性。

（3）\textit{不确定性引导}：UGOT 使用扩散深度估计模型 DiffDP 的去噪过程作为不确定性代理——在 $T$ 步去噪中，不确定性高的像素被更新的频率更高。不确定性量化为：

\begin{equation*}
U(z_{0:T} | I) = \frac{c(z_{0:T} | I) + M(c(z_{0:T} | M(I)))}{2}, \quad c(z_{0:T}) = \frac{1}{T} \sum_{t=T,\ldots,1} |z_t - z_{t-1}|,
\end{equation*}

其中 $M(\cdot)$ 为镜像函数。在 patch 内通过 softmax 归一化不确定性得到权重图 $W_{\mathrm{patch}}(i) = \mathrm{softmax}(-U_{\mathrm{patch}}(i))$，计算不确定性加权均值深度 $\bar{D}_{\mathrm{patch}}^U$ 和 $\bar{D}_{\mathrm{gt}}^U$，再求解不确定性引导的最优传输问题 $\mathrm{OT}_\epsilon(\bar{D}_{\mathrm{patch}}^U, \bar{D}_{\mathrm{gt}}^U)$，从而降低高不确定区域深度预测的影响。

（4）\textit{总体损失函数}：UGOT 将逐像素深度损失 $L_{\mathrm{dn}}$（来自 DNGaussian 的归一化 patch $L_2$ 损失）、最优传输深度损失 $L_{\mathrm{ot}}$ 和 RGB 渲染损失 $L_{\mathrm{color}}$ 组合为：

\begin{equation*}
\mathcal{L} = \lambda_1 L_{\mathrm{dn}} + \lambda_2 L_{\mathrm{ot}} + \lambda_3 L_{\mathrm{color}}.
\end{equation*}

\textbf{设计动机}：最优传输的核心优势在于其天然支持分布间的柔性匹配——它不要求渲染深度逐像素精确复现 MDE 预测，而是衡量将一个分布"搬运"到另一个分布的最小代价。这使得深度监督能够利用 patch 内的局部空间上下文信息，同时对 MDE 预测中的局部误差具有天然容忍度。不确定性引导则进一步确保在深度预测最不可靠的区域（如遮挡边界、无纹理区域）降低监督力度，避免错误先验误导优化。

In Depth We Trust（2026.04）~\cite{2604_05715} 系统研究了如何可靠利用单目深度先验，提出分离不适定几何来选择性地进行单目深度正则化，限制深度不准确性向已良好重建的 3D 结构传播。该方法强调从弱对齐的深度变体中学习的重要性，在不同 GS 变体和单目深度骨干上均展现出一致的几何精度和渲染质量提升。

在深度一致性方面，Cascade Pearson Correlation Loss（2025.05）~\cite{2505_22279} 提出分层深度引导泼溅（HDGS）框架，通过多尺度级联 Pearson 相关损失在多个空间尺度上强制渲染深度与估计深度先验之间的一致性。RDG-GS（2025.01）~\cite{2501_11102} 提出相对深度引导策略，利用相对深度而非绝对深度来精修高斯场，避免绝对深度的多视角不一致性。FatesGS（2025.01）~\cite{2501_04628} 利用单目深度排序信息监督块内深度分布一致性，并通过多视角投影特征优化深度绝对位置。

\subsubsection{几何约束与表面重建}

本路线介绍利用法线、表面连续性等几何约束提升 3DGS 几何精度与表面重建质量的方法。

几何约束旨在利用法线、表面连续性等额外信息来提升 3DGS 的几何精度，使其不仅能渲染出逼真的图像，还能恢复出准确的表面形状。本节首先详细介绍 Pi-GS 的稀疏视角表面重建管线，然后概述其他几何约束方法。

Pi-GS（2026.02）~\cite{2602_03327} 利用前馈网络 $\pi^3$（无参考点云估计网络，可直接从图像预测稠密深度图和相机参数）提供稠密初始化，并引入不确定性引导的深度监督、法线一致性损失和深度翘曲正则化来补偿几何不准确性。

\textbf{基线与问题}：稀疏视角下 COLMAP 的 SfM 流程产生的初始点云极度稀疏——例如在 Mip-NeRF 360 的 bicycle 场景中使用 24 个视角时，COLMAP 仅生成 1,028 个点，远不足以支撑 3DGS 的高质量优化。此外，PGSR（Planar-based Gaussian Splatting for Surface Reconstruction）等表面重建方法依赖多视角观测者裁剪（multi-view observer trim）来保证每个 3D 点被多个相机看到，但稀疏视角下该机制无法正常运作；opacity reset 操作也会在稀疏条件下丢失背景细节并引入伪影。

\textbf{核心改动}：Pi-GS 构建了一个完整的稀疏视角表面重建管线，包含四个关键环节——（1）利用前馈网络 $\pi^3$ 替代 SfM 提供稠密初始化点云（将点云规模提升约三个数量级）；（2）引入置信度感知的 Pearson 相关损失替代标准 $L_1/L_2$ 深度损失，实现结构一致性约束的同时避免对低置信度区域的过拟合；（3）设计网格边界掩码消除 $\pi^3$ 分块处理引入的法线伪影；（4）通过深度翘曲生成伪视角提供额外的多视角监督信号。

\textbf{实现原理}包含以下关键组件。

（1）\textit{稠密初始化}：Pi-GS 使用预训练的前馈网络 $\pi^3$（无参考点云估计网络）同时预测逐视角深度图和相机参数，替代 COLMAP 的 SfM 流程。在 bicycle 场景中，$\pi^3$ 在默认 20\% 置信度阈值过滤后仍产出 1,013,106 个点，相比 COLMAP 的 1,028 个点提升了约三个数量级。同时，Pi-GS 关闭了 PGSR 的多视角观测者裁剪（因稀疏视角下无法保证每个点被多视角看到）、关闭了 opacity reset（避免背景细节丢失）和分裂策略（因稠密初始化已提供充分的高斯覆盖）。

（2）\textit{置信度感知 Pearson 深度损失}：直觉上来说，$\pi^3$ 预测的深度图虽然在大致结构上可靠，但具体的数值精度有限，因此我们不应要求渲染深度逐像素精确匹配预测深度，而应只要求两者的分布形状一致——这正是 Pearson 相关系数所衡量的。标准 $L_1/L_2$ 深度损失容易导致模型过拟合于 $\pi^3$ 预测的有限精度，DNGaussian 的全局-局部深度归一化也因 $\pi^3$ 预测本身具有尺度一致性而变得不必要。Pi-GS 采用 Pearson 相关系数作为深度损失的基础度量，并融入 $\pi^3$ 提供的逐像素置信度 $C_i$：

\begin{equation*}
\mu_p = \frac{\sum_{i=1}^{N} C_i D^p_i}{\sum_{i=1}^{N} C_i}, \quad \mu_t = \frac{\sum_{i=1}^{N} C_i D^t_i}{\sum_{i=1}^{N} C_i}, \quad \hat{D}_p = D_p - \mu_p, \quad \hat{D}_t = D_t - \mu_t,
\end{equation*}
\begin{equation*}
P_{\mathrm{conf}} = \frac{\sum_{i=1}^{N} C_i \hat{D}^p_i \hat{D}^t_i}{\sqrt{(\sum_{i=1}^{N} C_i (\hat{D}^p_i)^2)(\sum_{i=1}^{N} C_i (\hat{D}^t_i)^2)}}, \quad \mathcal{L}_{\mathrm{pearson}} = 1 - P_{\mathrm{conf}},
\end{equation*}

其中 $D^p_i$ 为渲染深度，$D^t_i$ 为 $\pi^3$ 估计的深度。Pearson 相关系数衡量线性结构一致性而非逐像素精确匹配，能够恢复 $\pi^3$ 深度中缺失的高频细节，同时置信度加权确保低可靠区域的深度预测对优化影响较小。

（3）\textit{无伪影法线监督}：从深度图通过像素偏导数 $\partial z / \partial x$ 和 $\partial z / \partial y$ 计算法线。由于 $\pi^3$ 以 $14 \times 14$ 像素 patch 为单位处理图像，相邻 patch 的梯度在边界处不连续，导致法线图中出现网格状伪影。Pi-GS 构建一个 $14 \times 14$ 网格掩码，遮挡每个网格单元内 1 像素宽的内部边界区域，在损失计算时跳过这些不连续位置。法线损失采用 $L_1$ 范数：

\begin{equation*}
\mathcal{L}_{\mathrm{normal}} = \frac{1}{N} \sum_{i=1}^{N} \| N^t_i - N^p_i \|_1,
\end{equation*}

其中 $N^t_i$ 和 $N^p_i$ 分别为 $\pi^3$ 估计和渲染的法线。

（4）\textit{深度翘曲伪视角生成}：为进一步增强多视角监督，Pi-GS 通过深度翘曲在相邻相机之间插值生成伪视角。具体地，将源视角的图像像素反投影到 3D 空间，再重投影到目标相机平面。仅保留高置信度像素的投影结果，遮蔽不可见区域和低置信度区域。伪相机位置通过圆形插值确定——以目标相机和两个最近邻相机定义一个圆，在相邻视角对之间插值生成两个额外视角。伪视角通过 SSIM 和 $L_1$ 损失提供额外监督，权重设为 0.1。

\textbf{设计动机}：$\pi^3$ 网络的稠密初始化从根本上解决了稀疏视角下 SfM 失败的问题，而置信度感知的 Pearson 损失在利用结构一致性先验的同时避免了对预测误差的过拟合。网格掩码和深度翘曲则分别针对 $\pi^3$ 的特定缺陷（分块处理伪影）和稀疏视角固有的监督不足问题提供了针对性解决方案。

GTLR-GS（2026.03）~\cite{2603_23192} 将 LiDAR 数据作为主动几何先验融入整个高斯优化过程，提出几何-纹理感知分配策略和曲率自适应精修机制。DensifyBeforehand（2025.11）~\cite{2511_19294} 提出在优化之前进行致密化的策略，结合稀疏 LiDAR（Light Detection And Ranging，激光雷达，通过发射激光脉冲测量距离的主动传感器）数据和单目深度估计生成稠密点云，通过 ROI（Region of Interest，感兴趣区域）感知采样优先关注语义和几何重要区域。该方法绕过了标准 3DGS 中可能引入冗余高斯的自适应密度控制，使优化过程专注于高斯基元的其他属性（颜色、不透明度等），在保持与当时最优方法可比的渲染质量的同时显著降低了资源消耗和训练时间。

UNG-GS（2025.03）~\cite{2503_11172} 在 3DGS 管道中引入显式空间不确定性场（SUF）来量化几何不确定性，实现不确定性感知的深度渲染和法线精修。Reprojection-based DoF Separation（2024.12）~\cite{2412_14568} 提出将高斯位置自由度分离为图像平面平行 DoF（高置信度）和射线方向 DoF（低置信度），分别施加不同约束以抑制几何伪影。DET-GS（2025.08）~\cite{2508_04099} 提出统一的深度和边缘感知正则化框架，结合层次化几何深度监督、边缘感知深度正则化和 RGB 引导的边缘保留全变差损失。DIP-GS（2025.08）~\cite{2508_07372} 利用深度图像先验（Deep Image Prior，一种无需预训练的CNN结构本身即可作为自然图像先验的现象）的内部结构模式进行粗到细的 3DGS 优化，无需任何预训练模型。Few-shot Depth Aware GS（2024.10）~\cite{2410_11080} 使用单目深度预测作为先验配合尺度不变深度损失，并采用低阶球谐函数建模颜色以避免过拟合。

Depth Estimation Based on 3DGS（2024.09）~\cite{2409_12323} 利用 3DGS 的可微渲染来辅助单目深度估计的自监督训练。Enhancing Sparse-view（2025.01）~\cite{2501_11508} 通过 DINO-ViT 特征实现多视角语义一致性，并提出局部深度正则化约束深度值。Multimodal-Prior-Guided（2026.03）~\cite{2603_02866} 融合光度渲染残差、语义先验和几何先验来驱动分层高斯致密化。GauSSmart（2025.10）~\cite{2510_14270} 桥接 2D 基础模型和 3DGS，利用 DINO 语义特征引导高斯致密化和精修。

\begin{takeaway}
位姿估计与几何先验引导是稀疏视角 3DGS 的关键支撑：联合位姿-外观优化、前馈位姿估计和深度正则化三条路线互补，单目深度先验已成为弥补稀疏视角监督不足的标准组件，而 DUSt3R 等几何基础模型正在颠覆传统 SfM 管道。
\end{takeaway}

% ======================================================================


\section{退化处理与鲁棒重建}
\label{sec:12_deblur}

实际拍摄中，运动模糊、散焦模糊、低光照、噪声、雾、雨、烟等图像退化现象普遍存在，严重破坏 3DGS 所依赖的多视角一致性和 SfM 初始化。此外，真实世界场景中普遍存在瞬态干扰物（行人、车辆等）和复杂环境因素（水下散射介质等），进一步加剧重建困难。本节按退化类型和鲁棒性需求组织，分别讨论运动模糊处理、低光照增强、大气退化处理、多模态辅助去模糊、瞬态干扰物去除、水下与环境鲁棒重建，以及多模态增强与触觉融合。

\subsection{退化建模与恢复}
\label{ss:12_3a}

本节讨论四条技术路线：（1）运动模糊处理；（2）低光照增强；（3）大气退化处理；（4）多模态辅助去模糊。

\subsubsection{运动模糊处理}

本路线介绍从运动模糊图像中恢复清晰 3DGS 表示的去模糊方法。

运动模糊是最常见的退化类型，其核心挑战在于：模糊图像导致 SfM（运动恢复结构，Structure from Motion）位姿估计不准确，进而使 3DGS 初始化偏移。此外，模糊还导致 SfM 提取的特征点大幅减少，生成的初始点云极为稀疏，进一步恶化 3DGS 的起点。直觉上来说，运动模糊的本质是曝光时间内相机与场景的相对运动导致多个瞬时图像的叠加，因此去模糊的关键在于从模糊图像中分离出每一时刻的清晰状态。这一过程的困难在于模糊核通常是未知的，且不同像素位置的模糊程度可能不同。

Deblurring 3D Gaussian Splatting（2024.01）~\cite{2401_00834} 是该方向的开创性工作。在 Baseline 层面，标准 3DGS 直接用模糊图像训练时，由于模糊本质上是真实图像与点扩散函数（PSF，Point Spread Function，描述成像系统对点光源的模糊响应）的卷积，而 3DGS 无法建模这种邻域像素间的干扰，导致渲染结果同样模糊不清。此前基于 NeRF 的去模糊方法（如 Deblur-NeRF、DP-NeRF）通过体渲染中的逐像素卷积核来建模模糊，但这些方法依赖体渲染管线，无法迁移到基于光栅化的 3DGS 框架。

该工作的核心设计动机是：与其在图像空间中做反卷积去模糊，不如直接操纵 3D 高斯的几何属性来正向建模模糊过程。大尺寸高斯覆盖更宽的图像区域，恰好对应邻域像素对中心像素的干扰，这正是模糊的物理根源。具体而言，该方法引入一个小型 MLP $F_\theta$，以第 $j$ 个高斯的位置 $x_j$（经位置编码增强后输入）、四元数 $r_j$、缩放 $s_j$ 和视角方向 $v$ 为输入，针对两种模糊类型分别输出操纵参数：

\begin{itemize}
\item \textbf{散焦模糊建模}：MLP 输出缩放因子 $\Delta r_j$ 和 $\Delta s_j$，经截断和逐元素乘法后得到扩展的四元数 $\tilde{r}_j$ 和缩放 $\tilde{s}_j$。具体地，缩放因子经 $\sigma_s$ 缩放和 $(1-\sigma_s)$ 偏移后截断下界为 1.0：

\begin{equation}
\tilde{r}_j = r_j \cdot \min(1.0,\, \sigma_s \cdot \Delta r_j + (1-\sigma_s)), \quad \tilde{s}_j = s_j \cdot \min(1.0,\, \sigma_s \cdot \Delta s_j + (1-\sigma_s))
\end{equation}

其中 $\sigma_s = 0.01$ 控制缩放幅度。每个高斯的协方差被放大后承担更宽区域的信息，等价于在图像空间中模拟散焦扩散。推理时无需激活 MLP，直接用原始 $G(x_j, r_j, s_j)$ 渲染清晰图像，因此完全保留 3DGS 的实时性。
\item \textbf{运动模糊建模}：MLP 输出 $M$ 组偏移量 $\{\Delta x_j^{(i)}, \Delta r_j^{(i)}, \Delta s_j^{(i)}\}_{i=1}^{M}$，其中 $\Delta x_j^{(i)}$ 是第 $i$ 个时刻的位置偏移，$\Delta r_j^{(i)}$ 和 $\Delta s_j^{(i)}$ 是该时刻的旋转和缩放因子。位移偏移量经 $\sigma_p = 0.01$ 缩放后加到原始位置上：$\hat{x}_j^{(i)} = x_j + \sigma_p \cdot \Delta x_j^{(i)}$。由此构造出 $M$ 组高斯集合，光栅化生成 $M$ 幅清晰图像后取平均得到模糊图像 $I_b = \frac{1}{M}\sum_{i=1}^{M} I_i$。实验中设 $M=5$，即用 5 个离散时刻近似曝光时间内的相机运动。
\end{itemize}

MLP 本身采用 4 层结构：前 3 层为共享层（每层 64 个隐藏单元，ReLU 激活，Xavier 初始化），第 4 层分为 3 个独立头分别预测 $\Delta x$、$\Delta r$ 和 $\Delta s$。针对模糊导致的稀疏点云问题，该方法还提出从均匀分布中采样额外点并用 KNN 插值着色（最多补充 200,000 个点），以及基于深度的差异化剪枝策略（远处高斯使用更低的剪枝阈值以保留远端结构）。在 Deblur-NeRF 基准上，该方法在 SSIM 和 LPIPS 指标上达到当时最优，同时保持实时渲染速度。

DeblurGS（2024.04）~\cite{2404_11358} 进一步处理了由严重模糊导致的 SfM 位姿不准问题，估计每张模糊观测的 6-DoF（六自由度，即三维旋转加三维平移）相机运动来合成对应的模糊渲染，并提出高斯致密化退火（annealing，从高温度逐步降温的策略，此处指从抑制致密化逐步放开）策略防止训练早期生成不准确的高斯。该退火策略在训练初期通过抑制致密化操作来避免在位姿尚不精确时于错误位置生成高斯，随着相机运动估计逐渐收敛再逐步放开致密化，从而在真实世界和合成基准数据集以及手机拍摄视频上均达到当时最优的去模糊和新视角合成性能。

在连续运动建模方面，CRiM-GS（2024.07）~\cite{2407_03923} 利用神经常微分方程（ODE）预测连续相机轨迹，通过刚体变换和自适应畸变感知变换分别建模规则运动和滚动快门等非线性畸变。该方法通过回顾基础相机理论并结合神经训练技术，实现了对连续相机轨迹的精确建模——将曝光时间内的离散采样替换为连续积分，在运动模糊图像上恢复精细的 3D 场景同时保持实时渲染速度。CoMoGaussian（2025.03）~\cite{2503_05332} 在此基础上引入连续运动精修（CMR）变换，通过额外可学习参数细化刚体变换以更好地近似连续运动模糊的本质。

针对极端模糊，GeMS（2025.08）~\cite{2508_14682} 提出首个直接从严重模糊图像重建 3DGS 的框架，集成深度学习 SfM（VGGSfM）和 3DGS-MCMC 来避免传统 SfM 在严重模糊下的失败。BSGS（2025.10）~\cite{2510_12493} 提出双阶段策略：先精修位姿减轻运动畸变，再通过全局刚体变换进一步校正模糊畸变。BARD-GS（2025.03）~\cite{2503_15835} 将动态场景中的模糊分解为相机运动模糊和物体运动模糊两类，分别建模处理——前者由相机在曝光时间内的刚体运动引起，后者由场景中物体的独立运动引起。

CoherentGS（2025.12）~\cite{2512_10369} 专注于稀疏视角加运动模糊的复合退化，提出双先验策略：专用去模糊网络恢复清晰细节并提供光度引导，扩散模型提供几何先验来填充未观测区域。Casual3DHDR（2025.04）~\cite{2504_17728} 从随意拍摄的自动曝光视频中重建 HDR 3D 场景，将连续时间相机轨迹融入统一物理成像模型，联合优化曝光时间、轨迹和相机响应函数。

\subsubsection{低光照增强}

本路线介绍低光照条件下 3DGS 重建的增强方法。

低光照条件下的 3DGS 重建面临多视角亮度不一致、噪声放大和颜色畸变三重挑战。直觉上来说，低光照场景中不同视角往往使用不同的曝光参数以保持画面可用，这导致同一场景在不同照片中呈现截然不同的亮度和色调，给多视角一致性建模带来巨大困难。此外，低光照下的高 ISO 设置会放大传感器噪声，进一步降低图像质量。

Gaussian-DK（2024.08）~\cite{2408_09130} 是该方向的先驱工作。其核心洞察在于：暗光环境本身在物理世界中是均匀一致的，观测图像间的多视角不一致性本质上来自相机成像过程——不同视角采用不同的曝光参数（ISO、快门速度、光圈），导致同一场景在不同照片中呈现截然不同的亮度。不同视角采用不同曝光参数以确保主体清晰，导致图像亮度差异悬殊，而 vanilla 3DGS 无法处理这种不一致性。

基于此洞察，Gaussian-DK 对 3DGS 的表示和渲染管线进行了系统性改造。

\textbf{表示层面}，将球谐系数（Spherical Harmonics，用于表示视角依赖颜色的数学基函数）的语义从像素值（0--255）改为物理辐射值（0--$+\infty$），并全部在对数域中进行优化，以压缩值域和变化幅度，避免球谐函数在大幅值变化下的数值不稳定。

\textbf{渲染层面}，设计相机响应模块（Camera Response Module）$Z = f(H, E, L)$ 作为 3DGS 的调制函数，将辐射图映射为像素值的同时补偿视角间的不一致。该模块由三个组件构成：（1）\emph{曝光条件}（Exposure Level Condition）：从图像 EXIF 元数据中提取曝光时间 $T$、光圈 $A$ 和 ISO 感光度，计算经验曝光水平 $\mathcal{E} = T \cdot \text{ISO} / A^2$，作为全局亮度调制因子；（2）\emph{高斯光照特征}（Gaussian Light Feature）：为每个 3D 高斯附加一个可学习的光照特征向量，经泼溅生成 2D 光照特征图后，由轻量 CNN 学习光照特征与曝光水平之间的响应函数，输出逐像素亮度图 $M_{\text{light}} = \phi(\mathcal{E}, \mathcal{F})$，以捕捉复杂几何场景中不同位置对同一曝光水平的不同响应，弥补全局曝光条件在处理高频光照和阴影细节时的不足；（3）\emph{色调映射}（Tone-Mapper）：将调制后的辐射图 $\tilde{H} = H \cdot M_{\text{light}}$ 转换为像素值 $Z = g(\log(\tilde{H}))$，其中 $g$ 由逐通道（R/G/B）轻量 CNN 建模，以保留各通道色调映射的差异性，防止颜色退化。

此外，Gaussian-DK 观察到训练过程中少量高斯会试图补偿无法被相机响应模块完全建模的复杂亮度变化，它们在自适应密度控制中倾向于远离场景中心、靠近相机，形成"浮游物"伪影。为此提出基于训练迭代的梯度缩放策略，根据高斯在归一化设备坐标（NDC）中的深度计算距离权重 $w(i,j)$，在训练早期（$t < t_1$）对近相机高斯施加强梯度抑制以防止分裂和克隆，中期（$t_1 \leq t < t_2$）逐步退火，后期（$t \geq t_2$）完全释放以恢复高频细节。Gaussian-DK 还贡献了一个面向暗光场景的多视角基准数据集，并展示可通过控制曝光水平合成不同亮度的"照亮"图像。LE3D（2024.06）~\cite{2406_06216} 面向 RAW 图像的 HDR（High Dynamic Range，高动态范围，能够同时捕捉极亮和极暗细节的成像技术）视角合成，提出锥形散射初始化丰富 SfM 点云估计，用 Color MLP 替代球谐函数表示 RAW 线性色彩空间，并引入深度畸变和远近正则化改善场景结构。From Chaos to Clarity（2024.06）~\cite{2406_08300} 揭示了噪声对 3DGS 的特殊影响：3DGS 极易产生过拟合噪声的细长高斯形状，因此提出自监督框架集成噪声提取器和噪声鲁棒重建损失。

在 Retinex 理论框架下，LL-Gaussian（2025.04）~\cite{2504_10331} 引入端到端低光高斯初始化模块、双分支高斯分解模型（解耦反射率——物体表面固有颜色，和光照——场景中的照明条件）和无监督优化策略。Luminance-GS（2025.04）~\cite{2504_01503} 通过逐视角颜色矩阵映射和视角自适应曲线调整来处理多种挑战性光照条件（低光、过曝、变化曝光）。LLGS（2025.03）~\cite{2503_18640} 提出可分解高斯表示 M-Color，分别表征颜色信息进行定向增强。LL-GaussianMap（2026.01）~\cite{2601_15766} 首次将 2D 高斯泼溅引入低光图像增强，将增强任务建模为 2DGS 引导的增益图生成过程。LL-GaussianImage（2026.01）~\cite{2601_15772} 在 2DGS 压缩表示域内直接进行低光增强，避免了解压-增强-重压缩的繁琐流程。

LITA-GS（2025.04）~\cite{2504_00219} 提出光照无关的新视角合成方法，提取光照不变的物理先验来构建结构渲染策略和渐进去噪模块。DTGS（2025.11）~\cite{2511_13011} 首次将热成像监督引入低光 3DGS，通过循环增强-重建机制联合优化增强、几何和热监督。MERID-GS（2026.04）~\cite{2604_24053} 基于 Retinex 理论显式分离光照和反射率，通过可学习增益和光照状态引导的频率门控抑制噪声传播。FLED-GS（2026.03）~\cite{2603_08133} 将 3D 场景恢复重构为增强与重建的交替循环，插入中间亮度锚点实现渐进恢复。

针对 NTIRE 2026 挑战赛，一系列方法涌现：Naka-GS（2026.04）~\cite{2604_11142} 受生物视觉系统启发提出双分支 Naka 校正和渐进点剪枝。ELoG-GS（2026.04）~\cite{2604_12592} 集成学习式点云初始化和亮度引导的颜色增强。此外，DenoiseGS（2025.11）~\cite{2511_22939} 首次将 3DGS 应用于突发图像去噪，提出高斯自一致性损失和对数加权频率损失，在推理速度上比 NeRF 方法快 250 倍。DenoiseSplat（2026.03）~\cite{2603_09291} 构建了大规模噪声-清洁基准，在端到端前馈框架中实现噪声鲁棒的多视角重建。

\subsubsection{大气退化处理}

本路线介绍处理雾、烟、雨、大气湍流等大气退化条件下 3DGS 重建的方法。

大气退化（雾、烟、雨、大气湍流）通过散射和衰减严重破坏图像质量。本节按去雾、去雨、去烟雾和大气湍流四类退化分别介绍。

DehazeGS（2025.01）~\cite{2501_03659} 是首个基于 3DGS 的去雾方法。其核心创新在于直接在每个高斯基元上建立大气散射模型的传输函数，利用 3DGS 已显式建模的深度信息，省去 NeRF 方法中逐射线采样的开销。在 Baseline 层面，经典的大气散射模型（Atmospheric Scattering Model, ASM）将雾天观测建模为 $I(x) = J(x)t(x) + A(1-t(x))$，其中 $I$ 为观测雾图、$J$ 为潜在清晰图像、$A$ 为全局大气光、$t(x) = e^{-\beta d(x)}$ 为传输函数（$\beta$ 为散射系数，$d(x)$ 为场景深度）。此前基于 NeRF 的去雾方法（如 DehazeNeRF、ScatterNeRF）通过体渲染管线在射线级别建模传输，但体渲染的高采样成本严重制约了效率。DehazeGS 的核心设计动机是：3DGS 已显式建模了每个高斯的位置（即深度），因此可以直接在每个高斯基元上建立传输函数，省去 NeRF 方法中逐射线采样的开销，同时利用 3DGS 的多视角重建能力保证去雾结果的结构一致性。

具体而言，DehazeGS 对每个高斯基元 $G(x)$ 查询其深度 $d_G(x)$（归一化至 $[0,1]$ 并截断梯度以防止影响清晰高斯的潜在表示），并通过一个单层一维卷积网络 $F_\theta: d_G(x) \mapsto t_G(x)$ 学习深度到传输值的映射，其中 $\theta$ 包含可学习的散射系数。为保证输出在 $(0,1)$ 范围内，网络使用 ReLU 确保 $\beta \cdot d_G(x)$ 非负，再经 sigmoid 激活提升训练稳定性：

\begin{equation}
t_G(x) = \sigma\!\bigl(e^{-\max(0,\, \beta \cdot d_G(x))}\bigr)
\end{equation}

全局大气光建模为可学习参数 $A_G$，适用于所有高斯。雾天高斯的最终颜色通过大气散射模型直接作用在潜在清晰高斯上获得：

\begin{equation}
G_f(x) = G_c(x)\,t_G(x) + A_G\,(1 - t_G(x))
\end{equation}

其中 $G_c(x)$ 和 $G_f(x)$ 分别表示清晰和雾化的 3D 高斯。推理时只需移除传输项和大气光项，直接用 $G_c(x)$ 渲染即可得到去雾视角。

在训练目标方面，DehazeGS 提出多层次的损失函数组合。首先，渲染雾图与真实雾图之间的重建损失为标准 $L_1$ 与 D-SSIM 的加权组合：$\mathcal{L}_\text{rec} = (1-\lambda)\mathcal{L}_1 + \lambda\mathcal{L}_\text{D-SSIM}$（$\lambda=0.2$）。其次，通过 $\alpha$-混合 渲染传输图 $\mathcal{T}(P)$ 后，引入暗通道先验（DCP）和亮通道先验（BCP）作为物理引导损失。DCP 损失将传统暗通道先验重构为基于 3DGS 的能量函数 $\mathcal{L}_\text{DCP} = \mathbf{T}_D^T(P)\mathbf{L}\mathbf{T}_D(P) + \mathcal{T}^T(P)\mathcal{T}(P) - 2\mathbf{T}_D^T(P)\mathcal{T}(P)$，其中 $\mathbf{T}_D(P)$ 为 DCP 算法估计的传输图、$\mathbf{L}$ 为类拉普拉斯矩阵；BCP 损失 $\mathcal{L}_\text{BCP} = \|\mathbf{T}_B(P) - \mathcal{T}(P)\|_1$ 用于缓解 DCP 导致的暗化效应。此外，利用 DepthAnything V2 从雾图预测伪深度图作为深度监督 $\mathcal{L}_d = \|D_\text{pseudo} - D\|_1$，并在训练过程中对深度监督权重施加对数衰减（从 1.0 降至 0.01），以适应伪深度图的固有误差。最后引入深度加权重建损失 $\mathcal{L}_{d_\text{rec}} = \|d_\text{detach} \cdot (\hat{I} - I)\|_1$ 以增强远距离细节恢复。总损失为 $\mathcal{L} = \mathcal{L}_\text{rec} + \lambda_D\mathcal{L}_\text{DCP} + \lambda_B\mathcal{L}_\text{BCP} + \lambda_d\mathcal{L}_d + \lambda_{d_\text{rec}}\mathcal{L}_{d_\text{rec}}$。推理时无需激活传输网络，直接用原始清晰高斯 $G_c(x)$ 渲染，完全保留 3DGS 的实时性。

DeRainGS（2024.08）~\cite{2408_11540} 首次处理雨天环境下的 3D 重建，构建了 HydroViews 数据集。WeatherGS（2024.12）~\cite{2412_18862} 将天气伪影分为空中密集粒子（雪花/雨滴）和镜头遮挡两类，提出从密到疏的预处理策略依次移除。RaindropGS（2025.10）~\cite{2510_17719} 建立了镜头雨滴条件下的综合基准，揭示了聚焦位置对重建性能的影响。

在烟雾处理方面，SmokeSeer（2025.09）~\cite{2509_17329} 利用热成像和 RGB 图像双模态信息，通过 3DGS 融合两种模态并将场景分解为烟雾和非烟雾组件。GenSmoke-GS（2026.04）~\cite{2604_03039} 和 SmokeGS-R（2026.04）~\cite{2604_05301} 分别提出了基于物理引导的伪清洁监督和几何优先重建策略来处理烟雾退化。Dehaze-then-Splat（2026.04）~\cite{2604_13589} 采用两阶段管道：先用生成式去雾产生伪清洁图像，再用物理信息辅助损失训练 3DGS。Smoke-GS（2026.04）~\cite{2604_05687} 引入轻量级视角依赖的介质分支来捕获烟雾引起的方向性外观变化。

GSTurb（2026.02）~\cite{2602_22800} 处理大气湍流引起的像素位移和模糊，集成光流引导的倾斜校正和高斯泼溅建模的非等晕模糊。CoCoGaussian（2024.12）~\cite{2412_16028} 通过物理建模弥散圆（Circle of Confusion）来处理散焦模糊。RecGS（2024.07）~\cite{2407_10318} 迭代构建 3DGS 并在每轮通过低通滤波分离水下场景中的焦散。CompSplat（2026.02）~\cite{2602_09816} 显式建模压缩伪影的帧级特征，在严重压缩条件下保持重建质量。

\subsubsection{多模态辅助去模糊}

本路线介绍利用事件相机等多模态传感器辅助 3DGS 去模糊的方法。

事件相机因其高时间分辨率、高动态范围和无运动模糊的特性，成为辅助 3DGS 去模糊的理想传感器。与标准相机不同，事件相机不是以固定帧率输出完整图像，而是异步报告每个像素的亮度变化——这一特性使其能够以微秒级精度捕捉场景动态，为分解运动模糊提供了关键的时间信息。

E2GS（2024.06）~\cite{2406_14978} 是首个将事件数据融入高斯泼溅的方法。在 Baseline 层面，事件相机异步报告像素 $(x,y)$ 处的亮度变化事件 $e(x,y,t,p)$：当亮度变化超过阈值 $C$ 时触发，极性 $p \in \{+1, -1\}$ 指示变化方向，但不报告绝对强度值。此前基于 NeRF 的去模糊方法（如 E2NeRF）尝试融合事件数据，但受限于体渲染的计算开销。E2GS 的核心设计动机是：利用事件相机的高时间分辨率特性，在曝光时间内将模糊图像分解为多个时间步的清晰图像，并通过事件流提供帧间亮度变化的连续监督信号。

具体而言，E2GS 的流程分为预处理和训练两个阶段。在预处理阶段，给定一组模糊图像及其对应曝光时间内的事件流，首先将曝光时间等分为 $N$ 个时间间隔，得到 $N-1$ 个事件箱 $B_i = \{e_j(x_j, y_j, t_j, p_j)\}$。然后利用事件双积分（EDI，Event-based Double Integration，一种从事件流和模糊图像联合恢复清晰图像的方法）模型从模糊图像和事件流中恢复 $N$ 幅清晰图像：由于模糊图像是曝光时间内所有清晰图像的平均 $I_\text{blur} = \frac{1}{N}\sum_{i=1}^{N} I_i$，而相邻时刻的清晰图像之间满足 $I_{i+1} = I_i + \sum_j C p_j$（利用事件数据补偿亮度变化），由此可递推得到所有时间步的清晰图像集。这些去模糊图像被送入 COLMAP 进行 SfM，估计初始 3D 高斯坐标和 $N$ 个时间步的相机位姿——直接用模糊图像运行 COLMAP 往往会失败。

在训练阶段，E2GS 设计了两类互补的损失函数。第一类是\textbf{图像渲染损失}：利用每个视角的 $N$ 个位姿 $\{P_i\}_{i=1}^N$ 分别渲染 $N$ 幅图像 $\{\hat{I}_i\}_{i=1}^N$，取平均得到模拟模糊图像 $\hat{I}_\text{blur} = \frac{1}{N}\sum_{i=1}^N \hat{I}_i$，与真实模糊图像计算 $L_1$ 与 D-SSIM 的加权损失：$\mathcal{L}_\text{blur} = (1-w_\text{D-SSIM})\mathcal{L}_1 + w_\text{D-SSIM}\mathcal{L}_\text{D-SSIM}$。第二类是\textbf{事件渲染损失}：从 $N$ 幅渲染图像中随机选取两帧 $\hat{I}_n, \hat{I}_m$（$n < m$），转换为灰度后在对数域计算亮度差，除以阈值 $C$ 得到估计的事件计数图 $\hat{B}'_{nm}$，与真实事件箱 $B'_{nm}$ 计算均方误差 $\mathcal{L}_\text{event} = \|\hat{B}'_{nm} - B'_{nm}\|_2^2$。该损失利用事件数据的高时间分辨率特性，为模糊过程提供连续的亮度变化监督。最终损失为 $\mathcal{L} = \mathcal{L}_\text{blur} + w_\text{event}\mathcal{L}_\text{event}$，通过权重参数 $w_\text{event}$ 平衡两类监督信号。

EvaGaussians（2024.05）~\cite{2405_20224} 利用事件流显式建模运动模糊的形成过程，联合优化 3DGS 参数和曝光期间的运动轨迹。BeSplat（2024.12）~\cite{2412_19370} 首次从单张模糊图像和对应的事件流恢复清晰的高斯泼溅，通过 Bezier SE(3) 参数化联合学习场景表示和相机运动。该方法同时最小化合成模糊渲染与真实模糊图像之间以及合成事件与真实事件流之间的差异，是首个在单个模糊输入条件下实现高斯泼溅去模糊的工作。EaDeblur-GS（2024.07）~\cite{2407_13520} 引入自适应偏差估计器网络估计高斯中心偏差。

EBAD-Gaussian（2025.04）~\cite{2504_10012} 通过合成曝光期间的多个潜在清晰图像构建模糊损失函数，并利用事件流监督任意时刻的亮度变化。DiET-GS（2025.03）~\cite{2503_24210} 提出两阶段训练策略，结合事件双积分约束和扩散先验来同时实现准确的色彩和精细的边缘细节。EvFlow-GS（2026.04）~\cite{2604_22183} 利用光流提取事件边缘信息，提出端到端可学习的双重积分来联合优化事件流、位姿和 3DGS。DeblurSplat（2025.09）~\cite{2509_18898} 首次实现无需 SfM 的去模糊 3DGS，利用 DUSt3R 直接从模糊图像获取初始点云，并通过事件流解码潜在清晰图像提供监督。

Dark-EvGS（2025.07）~\cite{2507_11931} 是首个事件辅助的 3DGS 低光框架，提出三元组级监督（整体知识+颗粒细节+清晰渲染）和颜色调匹配模块。AsyncEvGS（2026.05）~\cite{2605_07192} 提出灵活的高分辨率异步 RGB-事件双摄系统，利用 VGGT 进行跨域位姿估计，并通过结构驱动的事件损失和视角特定一致性正则化实现稳定的高保真重建。SpikeNVS（2024.04）~\cite{2404_06710} 引入脉冲相机作为新型神经形态传感器，利用时域积分而非时域微分来提供锐利的场景表示。

此外，Robust Gaussian Splatting（2024.04）~\cite{2404_04211} 统一处理运动模糊、位姿不准和颜色不一致等多种退化，将运动模糊建模为相机位姿上的高斯分布。Photoreal Scene Reconstruction（2025.06）~\cite{2506_04444} 强调视觉惯性束调整（VIBA）和滚动快门物理成像模型的重要性，在 Project Aria 设备上实现 +2dB 的 PSNR 提升。GS-Blur（2024.10）~\cite{2410_23658} 利用 3DGS 构建了多样化运动模糊数据集用于训练去模糊网络。GaussFusion（2026.03）~\cite{2603_25053} 通过几何信息引导的视频到视频生成器来修复 3DGS 重建中的浮游物、闪烁和模糊伪影。RainyScape（2024.04）~\cite{2404_11401} 利用神经网络的光谱偏差特性先获取低频场景表示，再通过方向敏感梯度重建损失分离雨纹和场景细节。

除图像退化外，真实世界场景中普遍存在瞬态干扰物（行人、车辆等）和复杂环境因素（水下散射介质等），破坏了 3DGS 的多视角一致性假设。

\subsection{场景鲁棒性}
\label{ss:12_3b}

本节讨论三条技术路线：（1）瞬态干扰物去除；（2）水下与环境鲁棒重建；（3）多模态增强与触觉融合。

\subsubsection{瞬态干扰物去除}

本路线介绍识别并去除行人、车辆等瞬态干扰物以恢复多视角一致性的方法。

瞬态干扰物（如行人、车辆等仅在部分视角出现的物体）破坏了 3DGS 的多视角一致性假设，导致重建中出现浮游伪影。本节首先详细介绍代表性工作 SpotlessSplats 的完整技术方案，然后梳理后续的改进方法。

\paragraph{Baseline 与问题分析.}原始 3DGS 假设所有训练像素均来自静态场景，对行人、车辆等瞬态干扰物没有任何免疫力。已有鲁棒方法（如 RobustNeRF）仅依赖光度残差（photometric residuals）识别异常像素，容易过度拟合于颜色误差而非真正识别干扰物的语义身份。此外，3DGS 的逐图像 mini-batch 训练（每批次包含整张图像而非随机像素采样）使得基于分位数的截断估计器（trimmed estimator）假设——每个 mini-batch 含相同比例异常值——不再成立。

\paragraph{核心改动.}SpotlessSplats 将干扰物检测从光度残差驱动转向语义特征驱动：利用 Stable Diffusion 提取的自监督特征图，通过两种互补的聚类策略识别干扰物区域，并在 3DGS 优化过程中渐进式地将识别出的干扰物像素从训练中排除。同时，针对 3DGS 训练的特殊性，该方法还设计了三项关键的适配改进。

\paragraph{实现原理.}在预处理阶段，首先为每张输入图像 $\mathcal{I}_n$ 提取 Stable Diffusion 特征图 $\mathcal{F}_n$，形成像素级语义描述。干扰物检测提供两种策略：（1）\textbf{空间聚类}（spatial clustering）：对每张图像的特征图执行层次聚类（agglomerative clustering），将图像过分割为 $K{=}100$ 个空间连通区域，根据每个聚类内像素的重构残差计算内点概率 $P(c \in \mathcal{I})$，再将聚类标签传播回像素级掩码 $\hat{\mathcal{M}}_*(p)$；（2）\textbf{时空聚类}（spatio-temporal clustering）：训练一个 $1{\times}1$ 卷积 MLP 分类器 $\phi$，从像素特征预测逐像素内点概率 $\hat{\mathcal{M}}_* = \sigma(\phi(\mathbf{f}; \hat{\mathcal{M}}))$，分类器与 3DGS 模型交替优化。分类器损失为 $\mathcal{L}(\phi) = \mathcal{L}_\text{sup}(\phi) + \lambda_\text{reg}\mathcal{L}_\text{reg}(\phi)$，其中 $\mathcal{L}_\text{sup}$ 仅对高置信内点/外点像素施加监督（基于当前重构残差的阈值判断：$\tau_\text{inlier}{=}0.5$，$\tau_\text{outlier}{=}0.9$），对其余像素则完全依赖特征空间的语义相似性；$\mathcal{L}_\text{reg}$ 通过最小化 MLP 的 Lipschitz 常数正则化，确保相似特征映射到相似概率。

\paragraph{3DGS 适配设计.}（1）\textbf{预热调度}（scheduled sampling）：掩码并非立即生效，而是从伯努利分布 $\hat{\mathcal{M}}_p \sim \mathcal{B}(1 - (1{-}\hat{\mathcal{M}}_*)\cdot\gamma_t)$ 中采样，其中 $\gamma_t$ 为阶梯指数调度器（从 1 衰减至 0），保证早期仍有稀疏梯度回传至不确定区域，避免过早将干扰物区域完全屏蔽。（2）\textbf{直方图中位数估计}：为解决逐图像训练下异常值分布不均匀的问题，维护一个 $B$ 个桶的残差直方图（桶宽为渲染误差下界 $10^{-3}$），以折扣更新方式追踪残差分布，从中提取广义中位数 $\hat{m}$ 作为鲁棒估计基准。（3）\textbf{基于利用率的剪枝替代 opacity reset}：原 3DGS 周期性重置所有高斯不透明度的操作会严重干扰残差分布追踪，SpotlessSplats 改为追踪每个高斯投影位置的梯度利用率 $u_g = \frac{1}{|N_v(t)|}\sum_{t}\sum_{w,h}\|\nabla^{(t)}_{\boldsymbol{\mu}_g}\hat{C}^{(t)}_{h,w}\|_2^2$，当 $u_g < \tau_u$（$\tau_u{=}10^{-8}$）时剪枝，既实现了 floater 清理又避免干扰残差追踪。

\paragraph{设计动机与外观建模.}上述设计的核心动机在于：光度残差无法区分"真正的干扰物像素"与"尚未收敛的静态像素"，而语义特征天然地将同类物体（如所有行人像素）归为一组，提供更可靠的干扰物识别信号。预热调度和利用率剪枝则直接解决将鲁棒优化从 NeRF 的逐像素训练迁移到 3DGS 的逐图像训练时产生的技术困难。此外，该方法还引入了外观建模（appearance modeling），为每个相机视角优化一个 64 维 latent 向量，通过 MLP 映射为球谐系数的线性变换 $\hat{c}_i = a \odot c_i + b$，建模自动曝光和白平衡变化，防止 latent 过拟合为干扰物的逐图解释。

\paragraph{DeSplat: 体积渲染级干扰物分离}

DeSplat（2024.11）~\cite{2411_19756} 提出了一种优雅的解决方案：直接基于体积渲染将干扰物和静态场景元素分离——在每个相机视角中初始化用于重建视角特定干扰物的高斯，在 $\alpha$-混合阶段分别建模静态 3D 场景和干扰物。该方法无需外部语义信息或预训练模型，不牺牲渲染速度，并在三个干扰物基准数据集上达到与已有方法可比的效果。

\paragraph{PDF-GS: 渐进式多阶段过滤}

PDF-GS（2026.04）~\cite{2604_12580} 放大了 3DGS 固有的自过滤特性，通过渐进式多阶段优化逐步移除干扰物：先利用差异线索逐步过滤，再从净化后的高斯表示恢复精细细节。与 DeSplat 的显式干扰物建模不同，PDF-GS 依赖 3DGS 自身的过滤能力来渐进剔除瞬态元素。

\paragraph{DualSplat: 失败到先验框架}

DualSplat（2026.04）~\cite{2604_21631} 提出``失败到先验''框架：将首轮保守训练中的重建失败（瞬态物体通常表现为不完整碎片）转化为第二轮重建的显式先验，结合光度残差、特征不匹配和 SAM2 实例边界构建对象级伪掩码。该方法的核心洞察是瞬态物体仅出现在部分视角中，在保守初始训练中往往表现为不完整碎片，这些失败本身恰好可以转化为干扰物的显式先验；一个轻量级 MLP 在线精修伪掩码，逐渐从先验监督转向自一致性约束。

\paragraph{HarmoGS: 双层面冲突协调}

HarmoGS（2026.05）~\cite{2605_13073} 从图像空间监督和梯度级优化两个层面处理冲突：语义一致性引导的掩码学习像素级一致性分数来细化掩码，双视角冲突感知梯度协调策略将视角特定梯度旋转到正交配置以减少负面方向干扰。

\paragraph{Semantic-Guided 3DGS: 视觉语言模型类别感知去除}

在特定干扰物方面，Semantic-Guided 3DGS（2026.02）~\cite{2602_15516} 利用视觉语言模型进行类别感知的瞬态去除，通过累积 CLIP 相似度分数识别并剪枝干扰物高斯。

\paragraph{ForestSplats: 可变形瞬态场与超像素掩码}

ForestSplats（2025.03）~\cite{2503_06179} 提出可变形瞬态场和超像素感知掩码，无需视觉基础模型即可有效表示和分离瞬态元素，在内存效率上显著优于已有方法。

\paragraph{NexusSplats: 分层光照解耦}

NexusSplats（2024.11）~\cite{2411_14514} 通过分层光照解耦和结构感知的遮挡处理机制，在保持渲染质量的同时将参数量减少 65.4\%。

\paragraph{MW-GS: 微观-宏观投影增强}

MW-GS（2025.01）~\cite{2501_14231} 利用微观-宏观投影和小波采样来增强不受约束照片集的重建。

\paragraph{Sparse View Distractor-Free: 稀疏视角语义实体匹配}

Sparse View Distractor-Free（2026.03）~\cite{2603_01603} 专注于稀疏视角下的干扰物处理，利用 VGGT 进行语义实体匹配和视觉语言模型识别大静态区域。

\paragraph{CorrGS: 对应引导噪声鲁棒重建}

CorrGS（2025.01）~\cite{2501_14319} 提出对应引导的高斯泼溅用于噪声视频的鲁棒重建，通过渐进式测试时适应将噪声观测与清洁 3D 映射对齐。

\subsubsection{水下与环境鲁棒重建}

本路线介绍水下散射介质等特殊环境条件下 3DGS 的鲁棒重建方法。

水下场景重建面临光吸收、散射和浑浊度导致的颜色畸变、几何伪影和远处坍塌等多重挑战。与大气退化不同，水下的散射介质是持续存在的环境因素，需要在渲染管线中对其进行物理建模，而非简单地将其视为可去除的噪声。

OceanSplat（2026.01）~\cite{2601_04984} 是水下 3DGS 重建的代表性工作，针对散射介质引起的多视角不一致和几何伪影问题，提出了三目视角一致性约束、合成核线深度先验和深度感知不透明度调整三项核心改进。

\paragraph{Baseline 与问题分析.}标准 3DGS 的 $\alpha$-混合 渲染假设光线沿直线传播，但在水下散射介质中，光吸收和散射导致颜色畸变（偏蓝绿色）、远处物体细节坍塌，以及 3D 高斯沿射线过度分散（产生浮游伪影）。已有的双目立体一致性方法仅提供单一方向的几何约束，在散射环境中不够鲁棒。

\paragraph{核心改动.}OceanSplat 在物理水下图像形成模型的基础上，引入三目立体框架（水平+垂直正交基线）增强几何约束，推导基于对极几何的合成深度先验实现自监督深度正则化，并设计深度感知不透明度调整机制抑制散射方向上的浮游伪影。

\paragraph{实现原理.}（1）\textbf{水下渲染模型}：沿射线 $\mathbf{r}$ 通过 MLP $\theta_\text{med}$ 预测衰减 $\sigma^\text{attn}$（描述光线穿过水体的强度损失）、后向散射 $\sigma^\text{bs}$（描述水体向相机散射的环境光）和介质颜色 $c^\text{med}$（水体本身的颜色）。物体颜色 $C^\text{obj}$ 和介质颜色 $C^\text{med}$ 分开渲染：$C^\text{obj} = \sum_{i=1}^N T_i^\text{obj}\,\alpha_i\,c_i\,e^{-\sigma^\text{attn}z_i}$，$C^\text{med} = \sum_{i=1}^N T_i^\text{obj}\alpha_i\,\sigma^\text{bs}\,c^\text{med} + T_N^\text{obj}c_\text{med}e^{-\sigma^\text{bs}z_N}$，其中 $T_i^\text{obj}$ 为透射率。

（2）\textbf{三目视角一致性}：从相机姿态 $P_c$ 出发，沿水平和垂直方向平移生成虚拟视角 $P_h, P_v$（平移量 $b_h, b_v$），分别渲染图像 $I_h^\text{obj}, I_v^\text{obj}$，通过视差逆变形（disparity-based inverse warping）对齐至中心视角。视差图由 $d_h(x,y) = f_h b_h / D_c(x,y)$ 和 $d_v(x,y) = f_v b_v / D_c(x,y)$ 计算，其中 $f_h, f_v$ 为焦距，$D_c$ 为中心视角的渲染深度。总三目一致性损失为 $L_\text{tri} = L_\text{obj-stereo} + L_\text{full-stereo} + L_\text{smooth}$，其中 $L_\text{obj-stereo}$ 约束物体分量在虚拟视角间的一致性，$L_\text{full-stereo}$ 约束完整图像（物体+介质）与真实图像的对齐，$L_\text{smooth}$ 为边缘感知视差平滑正则。

\paragraph{合成核线深度先验与辅助机制.}（3）\textbf{合成核线深度先验} $D_\text{epi}$：从三目视角的视锥交集中采样高不透明度的高斯，投影到 $P_h, P_v$ 形成对应点对，通过对极几何三角化得到深度先验，以边缘感知 Log-L1 损失 $L_\text{epi}$ 正则化渲染深度。该先验完全自监督，无需外部深度传感器。（4）\textbf{深度残差损失}：$L_\text{res} = \frac{1}{N'}\sum_{i=1}^{N'}|D_c(x_i) - z_i|$，约束单个高斯的相机空间 $z$ 坐标与渲染深度对齐，防止高斯沿射线过度分散。（5）\textbf{深度感知不透明度调整}：在训练早期（$t < t_\text{adj}$），调整后不透明度 $\alpha'_i = (1{-}w)\alpha_i + w\alpha_i^d$，其中 $\alpha_i^d = f_\alpha(\alpha_i, z_i, \mathbf{d}_i)$ 由 MLP 根据深度和视角方向预测，在散射严重的方向上抑制高斯不透明度，促进次优高斯的剪枝；$t > t_\text{adj}$ 后 $w$ 衰减至零以消除推理开销。

\paragraph{设计动机.}垂直基线提供了水平基线无法捕获的正交几何约束——多基线立体在经典立体匹配中已被证明优于单基线方法。合成核线深度先验巧妙地利用了三目框架本身的几何关系，将三角化深度作为自监督信号，避免了对外部深度估计器的依赖。总训练目标为 $L_\text{total} = L_\text{photo} + \lambda_\text{tri}L_\text{tri} + \lambda_\text{epi}L_\text{epi} + \lambda_\text{res}L_\text{res}$，其中 $\lambda_\text{tri}{=}0.1$、$\lambda_\text{epi}$ 从 $0.4$ 退火至 $0.2$、$\lambda_\text{res}{=}0.01$。

\paragraph{UW-3DGS: 体素回归水下图像形成}

UW-3DGS（2025.08）~\cite{2508_06169} 提出可学习的基于体素回归的水下图像形成模块，建模空间变化的衰减和后向散射，并引入物理感知不确定性剪枝（PAUP）移除噪声浮游高斯。与 OceanSplat 的三目立体框架不同，UW-3DGS 侧重于通过体素化的图像形成模型来显式建模水下物理过程。

\paragraph{DualPhys-GS: 双路径衰减-散射优化}

DualPhys-GS（2025.08）~\cite{2508_09610} 通过双路径优化机制分别处理衰减和散射：RGB 引导的衰减优化模型结合 RGB 特征和深度信息，多尺度深度感知散射模型捕获不同尺度的散射效应。

\paragraph{RUSplatting: 解耦 RGB 通道水下学习}

RUSplatting（2025.05）~\cite{2505_15737} 提出基于水下衰减物理的解耦 RGB 通道学习，并引入帧插值策略配合自适应加权方案来增强稀疏视角一致性。

\paragraph{R-Splatting: 水下修复桥接与不确定性感知优化}

R-Splatting（2025.09）~\cite{2509_17789} 将水下图像修复与 3DGS 桥接，引入不确定性感知的不透明度优化将不透明度建模为随机函数以抑制光照变化引起的突变梯度响应。

\subsubsection{多模态增强与触觉融合}

本路线介绍利用触觉信号、跨传感器数据等非视觉信息增强 3DGS 重建的方法。

除视觉传感器外，触觉信号、跨传感器数据等非视觉信息也能为 3DGS 重建提供互补约束。本节简要介绍触觉增强和跨传感器融合方面的初步探索。

FusionSense（2024.10）~\cite{2410_08282} 和 Touch-Augmented GS（2025.08）~\cite{2508_07717} 探索了触觉信号对 3DGS 重建的增强作用。前者融合基础模型先验与视觉和触觉传感器的稀疏观测，后者引入空间选择性触觉测量来精修几何和外观，在遮挡严重场景中将 Chamfer Distance 降低超过 15 倍。No Calibration, No Depth, No Problem（2026.02）~\cite{2602_23559} 研究了跨传感器视角合成问题，提出匹配-致密化-整合方法来消除 RGB-X 传感器之间的标定需求。

\begin{takeaway}
退化处理与鲁棒重建将 3DGS 从理想实验室条件推向真实世界：运动模糊、低光照、大气退化等物理退化的显式建模，以及瞬态干扰物去除和跨模态融合，是 3DGS 走向实际部署的必经之路。
\end{takeaway}

% ======================================================================
\section{主动视觉与扩展应用}
\label{sec:12_active}

主动视觉与规划将 3DGS 从被动接收数据的表示方法提升为能够主动决定"下一步看哪里"的智能体。此外，超分辨率、多光谱重建、数据集生成等扩展应用也在不断拓宽 3DGS 在稀疏和退化条件下的适用边界。本节按不确定性量化方法、系统框架、特定应用、超分辨率、多光谱与特殊模态、数据集生成与应用工具等方向展开。

\subsection{主动重建}
\label{ss:12_4a}

本节讨论三条技术路线：（1）不确定性量化与信息增益；（2）主动重建系统框架；（3）特定应用的主动规划。

\subsubsection{不确定性量化与信息增益}

本路线介绍基于信息论量化 3DGS 重建不确定性以引导下一最佳视角选择的方法。

不确定性量化是 NBV（Next-Best-View，下一最佳视角）选择的基础。只有准确量化了"当前重建中哪些区域还不够可靠"，才能有针对性地选择信息增益最大的下一个观测视角。本节首先详细介绍基于 Shannon 互信息的 GauSS-MI 方法，然后梳理其他不确定性量化方案。

\textbf{GauSS-MI}（2025.04）~\cite{2504_21067}。

\emph{Baseline：}此前 NBV 选择中的信息增益通常依赖体积渲染中的射线-体素相交计数或启发式的覆盖率指标，缺乏对高斯基元自身可靠性的概率建模，导致无法在 3DGS 的泼溅渲染框架下严格量化``当前地图还剩多少不确定性可被新视角消除''。
\emph{改了什么：}直觉上来说，GauSS-MI 的核心思路是为每个高斯赋予一个"可靠度"评分，然后利用信息论工具计算"从某个新视角观察后，整体不确定性能降低多少"。具体而言，GauSS-MI 为每个 3D 高斯引入二值随机变量 $r^{[i]}$ 表示"该高斯对渲染是否可靠"，构建概率模型 $P(r^{[i]})\in(0,1)$，并在此基础上推导出 Shannon 互信息的闭式表达式：Gaussian Splatting Shannon Mutual Information（GauSS-MI）。这使得信息增益的计算可完全在 3DGS 光栅化管线中高效完成。
\emph{实现原理：}系统分两步运行。（1）\textbf{概率更新}：每当收到新观测 $Z_k$，计算渲染图像与观测图像之间的逐像素损失图 $L_k = \alpha_c\|C-\hat{C}_k\| + (1-\alpha_c)\|D-\hat{D}_k\|$，然后通过逆传感器模型将损失沿泼溅路径反向映射到每个高斯，以 log-odds 形式递推更新可靠性概率：$l_{1:k}^{[i]} = -\sigma_T T^{[i]}\sigma_L^{-1} L_k + l_{1:k-1}^{[i]}$，其中 $T^{[i]}$ 为高斯 $i$ 沿射线的累积透射率，$\sigma_T,\sigma_L > 0$ 为超参数。初始状态下所有高斯的可靠概率 $P_0(r^{[i]})=0.5$（无先验信息）。（2）\textbf{GauSS-MI 计算}：基于更新的概率模型，候选视角 $\mathbf{x}$ 的期望互信息推导为
\begin{equation*}
I(\mathbf{r};\mathbf{z}) = \sum_{j=1}^{n_z} P(z^{[j]}|M^{[j]}) \sum_{i}^{[j]} \bigl(-\log P(r^{[i]})\bigr)\,T^{[i]},
\end{equation*}

其中 $P(z^{[j]}|M^{[j]})$ 是基于像素亮度 $M^{[j]}$ 的测量先验（建模传感器噪声），内层求和直接在光栅化中与颜色同步渲染，无需额外射线投射。信息增益函数 $-\log P(r^{[i]})$ 确保可靠性低的高斯贡献更大的信息增益，与直觉一致。
\emph{设计动机：}（1）概率更新利用 3DGS 已有的渲染损失，无需引入额外网络；（2）log-odds 递推形式将概率更新转化为可光栅化的逐像素运算，实现高效性；（3）互信息推导基于标准 Shannon 信息论，保证理论严谨性。
\emph{关键参数：}逆传感器模型中 $\sigma_T$ 控制透射率对更新的调制强度、$\sigma_L$ 控制损失敏感度；损失图中 $\alpha_c$ 平衡颜色与深度权重；测量先验中亮度公式 $M = 0.299R + 0.587G + 0.114B$ 将 RGB 转换为亮度以建模噪声水平。

GAVIS（2026.05）~\cite{2605_30342} 提出各向异性可见性场，利用球谐函数表示每个粒子相对于训练视角的可见性，集成到贝叶斯网络不确定性感知的光栅化器中实现 200 FPS 的实时不确定性量化。

POp-GS（2025.03）~\cite{2503_07819} 从最优实验设计角度重新构建信息量化问题，提出 T-最优性和 D-最优性两种准则。FisherRF 系列方法（包括 Next Best Sense（2024.10）~\cite{2410_04680}）利用 Fisher 信息矩阵量化不确定性，并将其扩展到同时选择视角和触觉姿态。SA-ResGS（2026.01）~\cite{2601_03024} 提出自我增强残差 3DGS，通过三角化训练视角和外推渲染视角生成增强点云来估计场景覆盖度，并提出首个针对 3DGS 的残差学习策略。

\subsubsection{主动重建系统框架}

本路线介绍集成地图表示、视角选择和路径规划的完整主动重建系统。

完整的主动重建系统需要在 NBV 选择之外集成地图表示、路径规划和碰撞避免。直觉上来说，主动重建就像探索一个未知房间——你需要一个地图来记录已知区域（地图表示），一个策略来决定看向哪里（NBV 选择），以及一条安全路径到达目标位置（路径规划与碰撞避免）。

\textbf{ActiveGS}（2024.12）~\cite{2412_17769}。

\emph{Baseline：}此前主动重建系统要么使用 NeRF 等隐式表示（渲染速度慢、难以在线更新），要么使用纯体素或 TSDF 地图（缺乏高保真外观），且多数系统依赖耗时的射线投射来评估候选视角的信息增益，无法满足机器人实时规划需求。
\emph{改了什么：}ActiveGS 提出混合地图表示——将高斯泼溅地图（Gaussian Surfel，高保真重建）与粗粒度概率体素地图（空间占用建模与路径规划）融合——并引入高斯基元的置信度（confidence）建模，使系统能在同一框架中同时完成探索（探索未知区域）和利用（精修重建不足表面）两类任务。
\emph{实现原理：}系统以迭代式``建图--规划''循环运行，直到达到预定任务时间。%
（1）\textbf{混合地图表示}：体素地图 $V$ 以均匀划分的体素建模占用概率，用于碰撞检测和路径规划；GS 地图 $G$ 基于 Gaussian Surfel 表示，每个基元 $g_i$ 由位置 $\mu_i$、旋转 $q_i$、缩放 $s_i$、颜色 $c_i$、不透明度 $o_i$ 和置信度 $k_i$ 定义。

%
（2）\textbf{增量建图}：收到新 RGB-D 观测后，先渲染当前视角的颜色图、深度图和不透明度图，然后通过二值掩码 $B(p)=(O(p)<0.5)\wedge(\mathrm{avg}|I(p)-I^*(p)|>0.5)\wedge(D(p)-D^*(p)>\delta D^*(p))$ 识别需要增加新高斯的区域（不透明度低、颜色偏差大或出现新几何），将对应像素反投影为新的 Gaussian Surfel。训练损失为 $L = w_c L_c + w_d L_d + w_n L_n$，其中 $L_c$、$L_d$ 为颜色和深度的 $L_1$ 损失，$L_n$ 包含法向量余弦距离与全变差平滑项。每 5 步执行可见性检查，删除所有历史视角均不可见的高斯。%
（3）\textbf{置信度建模}：每个高斯的置信度 $k_i = \sigma_i \exp(\rho_i)$ 由两部分组成——距离加权的余弦相似度 $\sigma_i = \sum_{j \in S(g_i)}(1 - d_{ij}/d_{\mathrm{far}})\,\hat{v}_{ij}^{\top}n_i$ 衡量观测视角与高斯法向的对齐程度（近距和正面观测权重更高），方向离散度 $\rho_i = 1 - \|\bar{v}_i\|$ 衡量观测方向的多样性（$\rho_i \to 0$ 表示观测角度趋同）。高置信度意味着高斯被多角度充分观测，低置信度标识重建不足区域。%
（4）\textbf{视角效用}：候选视角 $v_i^c$（定义在自由体素中心的离散格点上）的综合效用为 $U_{\mathrm{total}}(v_i^c) = \lambda\,U_V(v_i^c) + U_G(v_i^c)$，其中 $U_V = N_u / |V|$ 为探索效用（可见未探索体素比例），$U_G = -\mathrm{mean}(K_i)$ 为利用效用（渲染置信度图的负均值，低置信度区域贡献更高利用价值），$\lambda$ 为探索权重。

最终 NBV 选择综合考虑效用与 $A^*$ 路径代价：$v^* = \arg\max_{v_i^c}\bigl(U_{\mathrm{total}}(v_i^c) - \beta\,U_{\mathrm{path}}(v_i^c)\bigr)$。
\emph{设计动机：}（1）混合地图分离了``空间理解''（体素地图，用于碰撞和未探索区域检测）与``外观重建''（GS 地图，用于高保真渲染），各取所长；（2）置信度直接从观测历史的空间分布推导，无需额外训练网络，且可通过 GS 渲染管线高效获取任意视角的置信度图；（3）视角采样结合随机局部采样与基于兴趣区域（ROI，前沿体素 + 低置信度高斯所在体素）的定向采样，兼顾全局覆盖和局部精修。
\emph{关键参数：}探索权重 $\lambda$ 平衡探索与利用；路径代价权重 $\beta$ 控制移动成本的影响；深度噪声常数 $\delta=0.05$；每步训练 10 次迭代（3 帧最近 + 5 帧随机）；候选视角数量 $N_{\mathrm{rand}} + N_{\mathrm{ROI}}$。

ActiveSplat（2024.10）~\cite{2410_21955} 建立了集成在线建图、视角选择和路径规划的统一框架，利用稀疏拓扑进行高效视角采样和路径规划，同时利用视角依赖的稠密预测进行视角选择。该方法的核心是混合地图表示——将环境的稠密信息与工作空间的稀疏抽象相结合，并采用基于拓扑图的层次化规划策略来减少重复轨迹、在有限时间预算下提升局部粒度，从而实现高保真重建和照片级新视角合成。

MAGICIAN（2026.03）~\cite{2603_22650} 突破了贪婪的最近最佳视角范式，提出长期规划框架：利用从预训练占用网络导出的``想象高斯''来高效计算任意新视角的覆盖增益，集成到树搜索算法中进行长期规划。该方法的``想象高斯''表示基于预训练占用网络的强结构先验，通过快速体渲染实现任意候选视角的覆盖增益计算，在闭环中迭代更新想象高斯并精修轨迹，在室内外基准上均达到当时最优性能。ActiveGAMER（2025.01）~\cite{2501_06897} 利用 3DGS 的高效渲染能力构建基于渲染的信息增益模块，结合从粗到精探索、后精修和全局-局部关键帧选择策略。GS-Planner（2024.05）~\cite{2405_10142} 和 HGS-Planner（2024.09）~\cite{2409_17624} 分别提出基于采样的主动重建策略和层次化规划框架。RT-GuIDE（2024.09）~\cite{2409_18122} 开发了 GPU 加速的运动规划算法，在保持重建质量的同时比已有方法快数个数量级。

$\mathbb{R}^3$-RECON（2026.01）~\cite{2601_07484} 提出了无需辐射场的主动重建框架，从轻量体素地图诱导隐式的姿态条件渲染性场，在毫秒级即可查询候选视角的渲染性分数。ActiveInitSplat（2025.03）~\cite{2503_06859} 专注于图像选择对高斯初始化的影响，基于密度和占据准则确保训练图像来自多样化视角。

\subsubsection{特定应用的主动规划}

本路线介绍美学摄影、物体级重建等特定应用场景下的主动规划方法。

在特定应用场景下，主动规划需要融入领域特定知识。本节按美学摄影、物体级重建和特殊场景三个类别组织相关工作。PhotoAgent（2026.03）~\cite{2603_22796} 利用大型多模态模型的推理能力将主观美学目标转化为可求解的几何约束，在 3DGS 构建的光真实感内部世界模型中进行迭代优化。Aesthetic Camera Viewpoint（2026.02）~\cite{2602_20363} 提出 3D 美学场概念，利用前馈 3DGS 网络将 2D 美学知识蒸馏到 3D 空间，实现稀疏捕获下的高效美学视角建议。SplaTraj（2024.10）~\cite{2410_06014} 将语言指令驱动的相机轨迹生成建模为连续时间轨迹优化问题。

在物体级重建方面，OUGS（2025.11）~\cite{2511_09397} 从高斯基元的物理参数（位置、缩放、旋转）直接推导不确定性，通过渲染 Jacobian 传播参数协方差，并集成语义分割掩码产生对象感知的不确定性分数。Informative Object-centric NBV（2026.02）~\cite{2602_08266} 提出实例感知的 NBV 策略，将实例级信息蒸馏为 one-hot 对象向量，计算置信度加权信息增益。

在特殊场景中，Perturbed Gaussian Ensemble（2026.03）~\cite{2603_06852} 针对稀疏视角 CT 重建中的 X 射线主动视角选择，通过随机密度缩放构建高斯密度场集成来建模不确定性。DAV-GSWT（2026.02）~\cite{2602_15355} 利用扩散先验和主动视角采样合成高保真的 Wang Tiles。Active Semantic Mapping（2026.01）~\cite{2601_12122} 集成 Octomap 和 3DGS 实现园艺环境的主动语义建图。Risk-Averse Path Planning（2025.10）~\cite{2510_06481} 将 NBV 选择与风险规避路径规划统一。Object-centered Data Acquisition（2026.04）~\cite{2604_19216} 实现手机端实时引导的物体中心数据采集。Camera Splatting（2025.09）~\cite{2509_15677} 将相机建模为 3D 高斯进行视角优化。IntelliCap（2025.08）~\cite{2508_13043} 利用语义分割和视觉语言模型引导用户采集。

Understanding while Exploring（2025.06）~\cite{2506_00225} 提出语义和几何不确定性量化引导的主动语义建图。UPNet（2025.06）~\cite{2506_14856} 利用轻量前馈网络预测不确定性图来加速视角选择。User-in-the-Loop（2025.06）~\cite{2506_21009} 通过局部重建的光场和误差峰值可视化引导用户采样。Fisher Information for Semantic and Dynamic（2025.12）~\cite{2512_22771} 将 Fisher 信息扩展到语义高斯参数和变形网络的主动学习。Frequency-based View Selection（2024.09）~\cite{2409_16470} 在频域排列候选视角来估计潜在信息增益。

\subsection{扩展应用}
\label{ss:12_4b}

本节讨论三条技术路线：（1）超分辨率；（2）多光谱与特殊模态；（3）数据集生成与应用工具。

\subsubsection{超分辨率}

本路线介绍从低分辨率输入训练的 3DGS 中渲染高分辨率图像的方法。

针对此前 3DGS 超分方法多采用单图像超分范式、直接对渲染低分辨率图像进行 2D 上采样，忽略了高斯表征中固有的多视角一致性约束、导致跨视角纹理不一致和棋盘伪影的局限，SRGS（2024.04）~\cite{2404_10318} 提出将 3DGS 超分辨率统一为先验注入和跨视角正则化的模块化框架。MVGSR（2025.12）~\cite{2512_15048} 首次引入核线约束的多视角注意力机制到 3DGS 超分辨率，通过基于相机姿态的辅助视角选择适配任意组织的多视角数据集。IE-SRGS（2025.11）~\cite{2511_22233} 联合利用外部 2D 超分辨率先验和内部 3DGS 特征，通过掩码引导融合策略协同利用两者互补优势。Close-up-GS（2025.03）~\cite{2503_09396} 通过渐进式自训练和 3D 感知生成模型增强近距视角的细节渲染。

\subsubsection{多光谱与特殊模态}

本路线介绍将 3DGS 扩展到热红外、近红外等多光谱领域的方法。

除可见光 RGB 之外，热红外、近红外等非可见光谱段也携带着丰富的场景信息。本节介绍将 3DGS 扩展到多光谱领域的方法。

针对传统多光谱重建依赖严格跨模态相机标定、标定误差直接导致重建几何错位；且热红外/近红外图像缺乏 RGB 纹理细节、直接优化会导致高斯颜色模糊的局限，MS-Splatting（2025.06）~\cite{2506_03407} 提出多光谱 3DGS 框架，无需跨模态相机标定即可从不同光谱域（热红外、近红外等）的图像生成多视角一致的新视角，利用神经颜色表示将多光谱信息编码为每个 splat 的紧凑特征嵌入。ThermalGaussian（2024.09）~\cite{2409_07200} 首次实现 RGB 和热红外双模态的 3DGS 重建，引入多模态正则化约束将存储成本降低 90\%。

\subsubsection{数据集生成与应用工具}

本路线介绍基于 3DGS 生成训练数据集和构建应用工具的方法。

针对现有图像匹配器依赖大量人工标注的稠密对应关系训练、数据集构建成本极高，且合成数据生成的对应关系缺乏几何精确性、导致在真实图像上泛化差的局限，MatchGS（2025.11）~\cite{2511_21265} 利用 3DGS 生成几何精确的稠密对应标签来训练零样本图像匹配器，将核线误差降低最多 40 倍。GS2E（2025.05）~\cite{2505_15287} 利用 3DGS 构建大规模合成事件数据集，通过自适应轨迹插值和物理一致的事件对比阈值建模生成时间密集且几何一致的事件流。Aerial-Ground Image Matching（2025.09）~\cite{2509_19898} 利用 3DGS 渲染中间视角来弥合航空和地面图像之间的巨大视角差异。NeRF-Based Augmentations for Spacecraft（2026.05）~\cite{2605_19649} 利用 NeRF 从少量图像生成多样化训练数据以训练航天器位姿估计器。Generating Synthetic Stereo（2025.06）~\cite{2506_04908} 利用 3DGS 生成立体数据集并探索知识迁移策略。PRIMU（2025.08）~\cite{2508_02443} 提出基于基元的不确定性估计框架，将训练误差和覆盖统计投影到高斯基元上进行后验不确定性预测。WE-GS（2024.06）~\cite{2406_02407} 提出高效的野外照片集 3DGS 框架，通过残差式球谐系数传递模块适应不同光照条件。GlowGS（2026.05）~\cite{2605_23602} 利用扩散模型和视觉基础模型补偿夜间场景缺失的结构线索。TUGS（2025.05）~\cite{2505_08811} 提出基于物理的水下场景紧凑张量化高斯表示。Rig-Aware 3D Reconstruction（2026.01）~\cite{2601_14208} 利用三摄相机组进行车辆底盘的自动 3D 重建。

\begin{takeaway}
主动视觉将 3DGS 从被动表示提升为主动决策智能体，而不确定性量化是实现"下一步看哪里"的核心；超分辨率、多光谱融合和数据集生成等扩展应用则持续拓宽 3DGS 在稀疏与退化条件下的适用边界。
\end{takeaway}

% ======================================================================

\section{近期研究（2026 年 8--11 月）}
\paragraph{DynActiveGS：分别衡量结构不确定性和动态干扰}
2026.08~\cite{2608_01178}。DynActiveGS 从动态环境的 RGB-D 观测估计像素不确定性，控制观测在高斯重建中的权重，再分别用结构和运动不确定性在 Voronoi 图上选视点。规划既要补足静态几何的覆盖，也要考虑动态物体可能造成的遮挡。

论文在 Social-MP3D 与动态增强的 Gibson 模拟场景检查重建与保留视图，保留视点不出现在各方法的训练轨迹里。结果对应模拟动态环境，传感器噪声和未建模动态仍需实机验证。

\paragraph{TRACE：用连续轨迹覆盖高斯地图的信息分布}
2026.08~\cite{2608_02304}。TRACE 从在线 2D 高斯面片地图聚合体素信息与原语不确定性，再优化可执行的连续遍历轨迹。视场已经覆盖的信息随轨迹规划逐步扣减，以减少对同一位置的重复访问。

在八个 Replica 场景中相对 ActiveGS 报告平均 PSNR 改善 1.5 dB。该值适用于所比较控制规划和重建协议，信息场质量仍取决于在线地图与不确定性估计。

\paragraph{CasDeblurGS：让两幅模糊图像经二维再到三维互相约束}
2026.08~\cite{2608_10345}。相机运动使图像模糊，而且从两幅图恢复视图时，直接跨图拷贝容易将遮挡和误匹配带入结果。CasDeblurGS 先用双向光流和一致性掩码得到局部二维指引，再由冻结的无位姿高斯骨干建立临时三维场景，以重渲染输入视角提供更密集的全局指引。修复后的两幅图交给同一骨干生成可渲染高斯。

作者在 Deblur-NeRF 固定 85 个两输入一目标的三元组，真实与合成子集相对最强对照的 PSNR 分别提高 1.19 和 2.11 dB。方法在测试时不需要输入位姿或逐场优化，但需要已知内参，研究设置为静态场景。极强模糊和大遮挡会降低跨视图引导的可靠性。

\paragraph{GaussVid：用带相机条件的视频扩散修复稀疏重建伪影}
2026.08~\cite{2608_21849}。GaussVid 通过受控高斯优化构造带伪影的成对视频数据，再用相机相关条件引导视频扩散修复视图。学习目标是让已有 3DGS 的跨视角渲染更一致，生成修复的图像细节仍可能缺少真实三维观测。

论文用 DL3DV 独立场景的 500 个片段评价，并用 Zip-NeRF 大型场景测试跨数据和重建方法迁移。应把画质修复、跨视角一致性与真实三维几何分别评价。

\paragraph{FujinSplat：在 RAW 域校正烟雾颜色后再拟合高斯}
2026.09~\cite{2609_06017}。FujinSplat 在 RAW 域学习紧凑可逆的颜色调整，控制器先在可给出精确标签的模拟 RAW 数据上训练，再对真实视图预测校正。高斯优化用有界的零均值颜色残差缓解不同视图间剩余的差异。

论文在 RealX3D 的六个烟雾场景对比，ConvIR 的平均 LPIPS 为 0.493，优于本方法的 0.541，而 PSNR 低 2.57 dB。全局颜色变换无法补回传感器已经丢失的纹理，也不代表可以直接迁移到任何相机的 RAW 流程。

\paragraph{TV-SGS：让三维原语以张量投票交换几何线索}
2026.09~\cite{2609_07734}。少视图下逐原语的光度约束不足以决定真实表面。TV-SGS 用张量投票在高斯间传播局部结构线索，再将得到的三维几何监督写入损失，优化时仍保持原有渲染任务。

DTU 的十五场景比较区分三视图的大重叠和小重叠输入。补充实验的较密输入下，带有多视图几何监督的 PGSR 骨干并未因投票项明显提高。因此应按照输入重叠和骨干原有几何监督理解收益。

\paragraph{Bi-FlowGS：让补全视频与高斯几何互相修正}
2026.09~\cite{2609_17039}。稀疏视图的扩散补全如果只提供伪颜色，场景可能用错误几何解释生成图。Bi-FlowGS 先从修复视频估计光流，令高斯渲染导出的运动对应与该流对齐；更新后的高斯流和几何再回馈视频修复，构成双向细化。生成数据提供跨视图线索，但其内容不自动成为实测结构。

在 Mip-NeRF 360 的 3、6、9 输入视图条件，原文相对各配置最强对照报告 PSNR 提升 0.43、0.34 和 0.38 dB。该基准没有稠密几何真值，部分几何指标来自共同的 COLMAP 轨迹。应区别渲染增益、与稀疏几何的一致性，以及没有条件确认的生成区域。

\paragraph{CADSplat：利用近似 CAD 对齐并约束少视图高斯}
2026.09~\cite{2609_18473}。CADSplat 检索一个近似 CAD 形状，用多视图轮廓对齐相机，并在高斯外观拟合时联合细化注册与非刚性形变。CAD 为不可见部分提供形状约束，图像仍决定可观测的外观和局部修正。

作者在 SCO-CAD、NeRS MVMC 和 NeRS Misc 上评价。主表的 SCO-CAD 保留视角结果含测试时位姿细化，部分基线未使用该细化。方法依赖合理的初始相对标定及近似可用的形状先验。

\paragraph{用适配后的前馈几何给稀疏水下视图补足匹配证据}
2026.09~\cite{2609_18737}。水下散射和吸收会同时污染颜色及几何伪标签。该方法保留清晰图的几何教师，用物理退化后的图训练 VGGT 学生，只给视觉聚合编码器装入 LoRA 并冻结预测头。适配后的 VGGT 从未知位姿输入预测稠密点，其高斯代理渲染相邻伪视图与深度，再让真实视图和伪视图共同参与 COLMAP，所得轨迹与深度用于初始化 RUSplatting。

原文在 SeaThru-NeRF 和 Submerged3D 各选四个场景，每八帧留出一帧测试。SeaThru-NeRF 上，整体流程将 COLMAP 初始化的 RUSplatting 平均 PSNR 从 24.37 提高到 27.11 dB。IUI3-RedSea 的强散射仍会破坏 SIFT 匹配，少视角下也有残余预测噪声，因此不能将两组数据上的结果外推到所有浑浊水域。

\paragraph{D3GS：用米制深度、跨视图特征和扩散辅助少视图优化}
2026.09~\cite{2609_22941}。D3GS 首先把来自 MapAnything 与束调整的稀疏米制深度补全后初始化高斯，再用深度、DINO 特征与扩散修复为优化提供互补监督。几何先验、跨视图对齐和生成外观分别约束不同误差来源。

论文报告少视图设定下的几何和画质比较。若初始米制深度失准，后续初始化可能继承误差；扩散生成补充的图像外观也需要通过跨视图观测判断是否对应真实几何。

\paragraph{GAPS：按信息价值和生成可靠度筛选伪视图}
2026.09~\cite{2609_23436}。GAPS 在少视图高斯的交替优化中，用深度、外观和场景适配约束生成器，再按候选视角的重建信息和生成可靠度安排补充，从保守内插逐步过渡到外推。通过检查的伪视图按像素不确定性进入训练，高密度区域另加高斯 dropout。

论文在 LLFF 和 Mip-NeRF 360 使用固定保留比例评价。LLFF 六视图相对原始 3DGS 平均 PSNR 高 0.89 dB。消融中移除视角选择会让 Mip-NeRF 360 的 LPIPS 劣于无伪视图的基线，生成先验的适用范围因此与筛选质量和场景覆盖有关。

\paragraph{LiTe-GS：用随机候选集减少主动选视的信息计算}
2026.09~\cite{2609_30393}。主动重建每次选下一视角时，如果对所有候选都计算 Fisher 信息，选视本身会成为开销。LiTe-GS 保留 POP-GS 的信息目标，在每轮贪心选择时只评估随机抽出的候选子集。近似参数控制样本量，选中的帧仍进入原有的逐场景高斯优化。

作者在 Blender 和 Mip-NeRF 360 上比较图像质量和 Fisher 评分调用次数。调用数下降而重建表现接近穷举选视的条件随子集大小而变。该效率结论度量选视阶段的调用数，不能直接视为端到端训练时间的同等降幅。

\paragraph{用双向轨迹可靠度调节渐进式无 COLMAP 重建}
2026.09~\cite{2609_30865}。该方法用 SE(3) 上的短时运动连续性提供相机位姿初值，再用正反向跟踪一致性衡量置信度，为渐进式重建中的传播和回溯调整提供控制信号。置信度用于筛选轨迹证据并减弱错误早期位姿向后传播。

论文把相机误差和保留视图画质分别列出，前者用 ATE 与相对位姿误差，后者用 PSNR、SSIM 等。双向一致性是内部可靠性线索，不能替代在急剧视点变化和跟踪失败时的外部位姿参考。

\paragraph{ClearGS：按可靠度保留手持视频的几何覆盖}
2026.09~\cite{2609_31509}。手持视频的模糊分布不均，直接删除低质量帧可能丢失某条轨迹上的全部观察。ClearGS 将视觉可靠度、退化风险和几何效用用于分配不同训练权重，保留弱覆盖区域的监督。修复以当前渲染作为位姿对齐的结构约束，以同一帧的原观测作为待恢复证据，最后汇总沿轨迹接受的修复。

论文在 GS2E 的三个模糊等级及 GSOTM 的模糊、滚动快门、两者结合和位姿噪声条件上测试。训练与修复过程没有配对清晰图，清晰目标只用于评价。感知质量改善不能证明修复得到的纹理都来自场景，作者也将修复证据的几何核对列为后续工作。

\paragraph{OC-GS：用固定轴运动模型校正不规则转台角度}
2026.09~\cite{2609_31572}。OC-GS 从不规则间隔的转台图像获得初始角度，在共享固定转轴的运动约束下联合细化角度和高斯几何。它利用的是可控的采集结构，所以能减少对一般自由运动相机估计的依赖。

在相同训练器和初始化的消融里，十二个不规则视图下细化角度比固定初始角度高 9.08 dB 前景 PSNR，六图时高 4.22 dB。该实验对应固定转轴扰动，方法仍依赖 VGGT 初始化和较小的转轴噪声。

\paragraph{AGILE-GS：用虚拟锚点和短名单减少主动视角筛选成本}
2026.09~\cite{2609_34176}。AGILE-GS 先在连续位姿空间优化虚拟锚点，以代表地图当前的信息方向，再用低冗余的候选短名单替代全候选渲染评估。基础模式避免逐候选 Fisher 重排，增强模式在短名单中加入重排。

作者在 Mip-NeRF 360、Tanks and Temples、Deep Blending 分开报告平均选点时间与重建画质，并另给具身采集结果。选点阶段的节省不等于相同比例的完整采集缩短，真实运动时间还由采集协议与可执行路径决定。

\paragraph{UGOD：用原语不确定性调节透明度和训练 dropout}
2026.09~\cite{2609_39089}。UGOD 从高斯属性与观察方向估计视图相关不确定性。它在透明度合成前连续抑制高不确定贡献，并以梯度分离的软 dropout 正则训练中的不可靠原语。前者作用于渲染，后者作用于优化。

论文在 LLFF 三视图及 Mip-NeRF 360 的八和二十四视图设定评价。LLFF 使用 COLMAP 重建的三个训练视图，不能把结果直接推断给未知相机参数或其他稀疏度。训练细节和画质结果应按协议对应。

\paragraph{MoonGS：从双视图和语义先验重建低纹理月面}
2026.10~\cite{2610_07110}。MoonGS 从一对带相机参数和语义图像的月面视图直接预测高斯。方法利用视觉基础模型的深度特征确定基元位置，在非重叠或缺少纹理的区域以语义先验改善深度估计，并用深度排序约束稳定背景。启发式视图重采样用于补足稀疏观测。

作者在 LuSNAR 和合成 MoonBlender 上以重训的前馈新视图模型为基线，LuSNAR 的 PSNR 由所比基线的最高 18.89 dB 提高到 23.80 dB。真实嫦娥图像只用于定性比较。该增益不能直接解读为真实月面的米制几何误差下降，两幅输入也限制了大面积遮挡区的证据。

\paragraph{位姿图替代束调整的失效分析}
2026.10~\cite{2610_14090}。作者检验密集闭环下用仅包含位姿的图优化替代全局束调整。实验中的子块仍需做局部束调整，跨块再用位姿约束合并。纯 SE(3) 边不能表达各子块尺度差异，成对边也不能恢复同一个三维点在多视图间形成的联合约束。

在两个详细展示的自采无人机场景及公开 GauU-Scene V2 中，所测实现未达到调优 COLMAP 精度。消费级 GPS、依序合并和手工权重也可能造成差距。由于没有对这些选择分别做消融，论文的负结论不能简单扩大到所有分布式估计或保留稀疏结构的混合方法。

\paragraph{以少样本训练控制动态稀疏视图高斯过拟合}
2026.10~\cite{2610_25887}。作者把动态稀疏视图重建的过拟合拆为监督不足、相机偏差和模型容量三个问题。光流用于合成额外监督视图，重投影误差用于联合调整相机与场景。训练端选择性丢弃静态高斯，并以早停减少对少量观测的共同适配。

定量表来自 SIGGRAPH Asia 2025 稀疏视角竞赛的两个验证序列，每个用八个训练视角与八个新视角评价，只计前景区域。伪视角生成假定物体运动适度，特别快的人体动作可能破坏该条件。

\paragraph{CrossFit-GS：用目标帧未参与训练的子图精调位姿}
2026.10~\cite{2610_28190}。逐帧推进的无 COLMAP 高斯重建可能把某帧的位姿误差吸收到高斯外观中，再用同一地图确认错误位姿。CrossFit-GS 暂缓把候选位姿写入全局地图，先以未使用目标帧且参数冻结的子图精化位姿，通过另留出的像素决定是否采用结果。

作者在 Tanks and Temples 与 CO3D-V2 同时评价轨迹和新视图。消融表明，对比只在原图上降低光度损失，隔离目标观测更有助于避免误差回流。交叉拟合仍需有足够重叠以建立可用参考，并要保守处理覆盖不足的更新。

\paragraph{DomeGS：用带视差的高斯天空穹顶完善 VR 远景}
2026.10~\cite{2610_37172}。DomeGS 从已有室外高斯场抽取全景，以考虑全周一致性的生成外扩补全背景，再把全景转成几何适配的多视角监督。背景最终作为天空穹顶高斯与原场景共同渲染，头部移动时可以产生经典无限远天空盒缺少的视差。

作者用 36 人的被试内 VR 研究比较天空盒与高斯穹顶，主要任务引导参与者观察地平线和上半球。结果支持测试场景下的背景一致性与临场感，补出的远景不是实际测量几何。

\paragraph{谱系感知稀疏高斯训练：阶段式随机屏蔽与均衡增密}
2026.11~\cite{2611_17752}。克隆及分裂可能产生协同拟合同一训练视角的一族高斯。该工作以增密谱系记录原语关系，早期对整个谱系共享随机屏蔽，随后改为独立屏蔽；同时在不改变克隆与分裂总配额的前提下，调整各谱系获得的增密机会。

论文在 LLFF、Mip-NeRF 360 与 Blender 中比较，并测试把两个模块接入不同主干的效果。作者报告收益因场景、指标和主干而异，阶段式屏蔽是主要贡献，均衡分配提供额外改进。

\paragraph{QuantileGS：为冻结高斯渲染器校准新视角误差上界}
2026.11~\cite{2611_21253}。QuantileGS 固定已有三维高斯的 RGB 渲染，再给原语拟合有序误差分位属性。留出的校准视角估计归一化校准量，从而为新视角给出随位置变化的 RGB 平均绝对误差上界。查询需要相机参数，不需要目标图像。

作者在 LF、LLFF 和 Mip-NeRF 360 比较误差排序与覆盖，并展示主动选视、裁剪和 SLAM 的经验收益。有限样本覆盖是在校准与测试误差分数可交换时对固定渲染器成立的边际保证，不能扩展为每个像素、每个视图同时有保证。


\section*{本章小结}
\addcontentsline{toc}{section}{本章小结}

\begin{futurebox}{本章总结与展望}{ch12_summary}
本章系统梳理了 3DGS 在稀疏条件和鲁棒重建方面的 268 篇研究论文，覆盖了从稀疏视角新视角合成、位姿估计与几何先验引导、退化处理与鲁棒重建到主动视觉与扩展应用的完整技术栈。通过分析可以发现以下核心趋势：

第一，\textbf{基础模型正成为稀疏视角 3DGS 的核心基础设施}。DUSt3R、MASt3R 等几何基础模型以及 Stable Video Diffusion 等生成模型，正在从根本上改变稀疏视角重建的范式——从依赖 SfM 管道的被动优化转向利用强先验的主动重建。

第二，\textbf{生成式先验与传统重建的融合}正成为解决不适定问题的主流路线。扩散模型不再仅用于后处理修复，而是作为训练监督的核心组成部分，通过 SDS、伪视角生成和视频补全等多种方式深度融入 3DGS 优化。

第三，\textbf{多模态感知}（事件相机、热成像、触觉、LiDAR）为退化条件下的鲁棒重建提供了互补信息通道，特别是在运动模糊、低光照和水下场景中展现了独特优势。

第四，\textbf{主动视觉}正在将 3DGS 从静态表示方法转变为动态的感知-决策闭环系统，使得机器人能够在未知环境中自主规划采集路径以最大化重建质量。

展望未来，稀疏条件下的 3DGS 重建仍面临若干开放挑战：如何在极少视角（1--2 张）下实现完整场景重建、如何有效融合异构传感器数据、如何在极端退化条件下保证重建鲁棒性，以及如何将主动规划与实时 3DGS 训练更紧密地集成。这些问题的解决将进一步推动 3DGS 在机器人、自动驾驶、数字孪生等领域的实际部署。
\end{futurebox}
